%% Supplementary material for The Visual Computer submission
%% "Fidelity Metrics Do Not Reliably Rank Visible-Infrared Fusion Rules for
%% Pedestrian Detection".
%%
%% Material moved here to keep the main text readable. Displays are numbered
%% S1, S2, ... and the shared bibliography is the one of the main text, so the
%% two lists agree. Build from _submission_IPT with pdflatex so that the figures
%% resolve through \graphicspath below.
\documentclass[10pt,a4paper]{article}
\usepackage[a4paper,margin=1in]{geometry}
\usepackage{mathptmx}
\usepackage{amsmath,amssymb,graphicx,booktabs,siunitx}
\usepackage{caption}
\usepackage{etoolbox}
\usepackage[numbers]{natbib}
\usepackage[hidelinks]{hyperref}

\graphicspath{{../}{./}}

\captionsetup{skip=6pt}
\AtBeginEnvironment{center}{\setlength{\topsep}{3pt}%
  \setlength{\partopsep}{0pt}}
\renewcommand{\thefigure}{S\arabic{figure}}
\renewcommand{\thetable}{S\arabic{table}}
\captionsetup[figure]{name=Fig.}
\captionsetup[table]{name=Table}

\begin{document}

\title{Supplementary Material for ``Fidelity Metrics Do Not Reliably Rank
Visible--Infrared Fusion Rules for Pedestrian Detection''}
\author{Jianyong Zhao, Zigao Wu, and Mengfei Ma}
\date{}
\maketitle

This supplement collects the material that supports the main text without being
needed to follow its argument: the remaining fidelity metrics for one benchmark,
the temperature sweep and its reading at matched weight deviation, the diagnostics
behind the classification-saliency comparison, the multi-scene qualitative
comparison, the definitions of the adaptation variants, a comparison of the
fidelity index under the reported index and the reference construction, the per-seed
detection contrasts, and
the segmentation results for the control family, the fidelity intervals for the
scrambling controls, the three-way rescoring of the fidelity index, and a protocol
audit of the detection comparison. Tables and figures are numbered
S1, S2 and so on, and the reference list is the one of the main text, so the two
lists agree.

\section{Additional fidelity metrics on RoadScene}

Table~\ref{tab:qtable} reports CC, SCD, VIF and SSIM on RoadScene for the same
eight methods. The two end-to-end state-of-the-art networks dominate VIF
(SPGFusion 0.83), while the classical rule and the targetness rules are
competitive on CC, SCD and SSIM.

\begin{center}\small

\captionof{table}{Quantitative comparison on RoadScene (means over 221 pairs).
Best in bold; the two end-to-end state-of-the-art networks dominate VIF, while the classical
baselines are competitive on CC/SCD/SSIM. SSIM is the mean of the two
single-reference SSIM values, against the visible and against the infrared source,
using the same convention as the main text.}\label{tab:qtable}
\begin{tabular*}{\linewidth}{@{\extracolsep{\fill}}l SSSS@{}}
\toprule
Method & {CC} & {SCD} & {VIF} & {SSIM}\\
\midrule
mean & \textbf{0.6662} & \textbf{1.3324} & 0.3127 & 0.7042\\
top-hat & 0.6487 & 1.2975 & 0.4241 & 0.6968\\
probe & 0.6566 & 1.3131 & 0.4227 & 0.6972\\
ensemble & 0.6544 & 1.3087 & 0.4234 & 0.6981\\
IFCNN~\cite{zhang2020} & 0.6452 & 1.2904 & 0.4988 & 0.6919\\
DenseFuse~\cite{li2019} & 0.6621 & 1.3242 & 0.3828 & \textbf{0.7146}\\
PIAFusion~\cite{tang2022} & 0.5931 & 1.1861 & 0.6316 & 0.6582\\
SPGFusion~\cite{zhang2026} & 0.6458 & 1.2915 & \textbf{0.8298} & 0.6455\\
\bottomrule
\end{tabular*}

\end{center}

\section{Sharpness control: temperature sweep}

The activity term of the rule also changes how sharply the softmax separates the
two sources, which is why the raw comparison in the main text mixes content with
sharpness. Table~\ref{tab:sharp} reports VIF and the mean absolute weight
deviation $\overline{|w-0.5|}$ as the temperature $\tau$ varies, and the two
paragraphs below read the rules at matched deviation instead of at matched
$\tau$.

At matched weight deviation the targetness maps are consistently better than the
activity-only rule rather than worse. On TNO, at a deviation of $0.037$, the probe
reaches VIF $0.3476$ and the top-hat prior $0.3381$ against $0.3087$ for the
activity-only rule, and at a deviation of $0.117$ the top-hat prior reaches
$0.4252$ against $0.3880$. On MSRS the activity-only rule is almost flat in
sharpness, moving only from $0.2911$ to $0.3102$ across a deviation range of
$0.018$ to $0.316$, while the top-hat prior reaches $0.4923$ at a deviation of
$0.092$; interpolated to the same deviation of $0.0525$ the prior is $+0.132$
above the activity-only rule. The targetness content therefore does carry
information, but that only becomes visible once the sharpness of the two
weightings is matched. Two configurations are close enough in deviation to be
compared pairwise rather than by interpolation. On TNO the top-hat prior at
$\tau=0.1$ (deviation $0.1255$) exceeds the activity-only rule at $\tau=0.3$
(deviation $0.1172$) by $+0.0454$ VIF with a paired interval of
$[+0.0423,+0.0486]$ over the 74 images, and on MSRS the prior at $\tau=0.1$
(deviation $0.0917$) exceeds the activity-only rule at $\tau=0.3$ (deviation
$0.1589$) by $+0.1914$ with $[+0.1832,+0.1998]$ over the 361 images. The MSRS
pair is not sharpness-matched, its two deviations differing by a factor of $1.7$,
so that value should be read as an upper bound; the TNO pair differs by
$7$ percent and is effectively matched.

\begin{table}[!t]\centering\small

\caption{Sharpness control. VIF and the mean absolute weight deviation
$\overline{|w-0.5|}$ as the softmax temperature $\tau$ varies. Reading the rules at
matched deviation instead of matched $\tau$ separates the content of the saliency
map from the sharpness of the weighting.}\label{tab:sharp}
\begin{tabular*}{\linewidth}{@{\extracolsep{\fill}}l l SSSS@{}}
\toprule
Rule & {$\tau$} & \multicolumn{2}{c}{TNO} & \multicolumn{2}{c}{MSRS}\\
\cmidrule(lr){3-4}\cmidrule(lr){5-6}
 & & {$\overline{|w-0.5|}$} & {VIF} & {$\overline{|w-0.5|}$} & {VIF}\\
\midrule
activity-only & 3 & 0.0124 & 0.2820 & 0.0177 & 0.2911\\
 & 1 & 0.0370 & 0.3087 & 0.0525 & 0.2923\\
 & 0.3 & 0.1172 & 0.3880 & 0.1589 & 0.3009\\
 & 0.1 & 0.2678 & 0.4835 & 0.3159 & 0.3102\\
\addlinespace
top-hat prior & 3 & 0.0049 & 0.2790 & 0.0035 & 0.3022\\
 & 1 & 0.0146 & 0.2993 & 0.0104 & 0.3244\\
 & 0.3 & 0.0474 & 0.3561 & 0.0340 & 0.3915\\
 & 0.1 & 0.1255 & 0.4334 & 0.0917 & 0.4923\\
\addlinespace
learned probe & 3 & 0.0026 & 0.2750 & 0.0020 & 0.3033\\
 & 1 & 0.0079 & 0.2883 & 0.0059 & 0.3282\\
 & 0.3 & 0.0258 & 0.3298 & 0.0189 & 0.4017\\
 & 0.1 & 0.0727 & 0.4044 & 0.0488 & 0.4980\\
\bottomrule
\end{tabular*}

\end{table}

The sharpness confound applies to the detection comparison as well, so we repeat
the temperature sweep for the three rules on the 300 LLVIP images and score every
setting with both task networks. The plain activity rule is fragile to sharpness.
Its mAP50 under the visible-trained detector falls from $0.8308$ at $\tau=3$ to
$0.8161$ at $\tau=1$ and to $0.7860$ at $\tau=0.1$ as the mean absolute weight
deviation grows from $0.022$ to $0.371$, whereas the top-hat prior and the probe
stay within $0.823$ to $0.831$ across the whole range. At the operating
temperature of the control experiment, $\tau=0.3$ and a weight deviation of
$0.196$ for the activity rule, the three rules are indistinguishable ($+0.0023$
and $-0.0027$, intervals crossing zero). At $\tau=0.1$ the two targetness rules
are significantly better than the activity rule, by $+0.0365$
$[+0.0064,+0.0698]$ for the prior and $+0.0440$ $[+0.0145,+0.0751]$ for the
probe. Under the modality-agnostic detector the targetness rules are at least as
good as the activity rule at every temperature and better at the two sharpest
settings, by up to $+0.0818$ $[+0.0569,+0.1082]$.

\section{Diagnostics of the classification-saliency rule}

\begin{sloppypar}
The classifier that the classification-saliency rule of Xu et
al.~\cite{xu2021} uses is a modality classifier, and the two diagnostics below
show why its saliency behaves like mean fusion on our benchmarks. We repeat the
channel-replacement diagnosis of Fig.~\ref{fig:csfch}, together with the
interpolation diagnosis of Fig.~\ref{fig:csfgrad}, on in-domain MSRS pairs.
Figs.~\ref{fig:csfch_a} and~\ref{fig:csfch_b} show the curves on six
representative pairs: a dip marks a channel whose replacement by the visible
feature makes the classifier lose confidence that the input is infrared, while a
channel that stays at the original probability carries redundant information.
Two properties are worth recording. First, the dips are concentrated in a few
channels and their position varies with the scene, so the per-channel importance
is scene-dependent rather than universal: over the whole MSRS test set the median
variation of the infrared probability is only about $0.02$ and roughly a third of
the scenes vary by more than $0.1$. Second, on scenes whose visible source is
nearly dark the classifier is saturated near $P(\mathrm{IR})\!\approx\!1$ and the
curve is flat, which is exactly the regime in which a modality-cue saliency
carries little information.
\end{sloppypar}

\begin{center}

\includegraphics[width=0.72\linewidth]{figures/csf_channel.png}
\captionof{figure}{Channel-replacement diagnosis of the classification-saliency
(CSF) classifier on an MSRS test pair. For each channel the infrared feature is
replaced by the visible feature and the infrared-class probability is re-evaluated;
dips (here channels 8, 10 and 23) mark channels that carry modality-specific
infrared information, whereas the flat remainder is modality-agnostic, which is why
the resulting classifier-saliency rule behaves like mean fusion.}\label{fig:csfch}

\end{center}

\begin{center}

\includegraphics[width=0.98\linewidth]{figures/diag_csf_channel_a.png}
\captionof{figure}{Channel-replacement importance, laid out as in Fig.~4
of~\cite{xu2021}, on MSRS
pairs 00357D, 00621D and 00790N. A dip marks a channel whose replacement by the
visible feature makes the classifier lose confidence that the input is infrared,
so the corresponding feature map carries modality-specific information.}\label{fig:csfch_a}

\end{center}

\begin{center}

\includegraphics[width=0.98\linewidth]{figures/diag_csf_channel_b.png}
\captionof{figure}{The same diagnosis on three further MSRS pairs (00738N, 00504D,
00462D). The dip pattern is scene-dependent but stays concentrated in a small number
of channels, and the flat segments are modality-agnostic.}\label{fig:csfch_b}

\end{center}

\begin{center}

\includegraphics[width=0.66\linewidth]{figures/diag_csf_gradient.png}
\captionof{figure}{Interpolation diagnosis, laid out as in Fig.~6
of~\cite{xu2021}. The input
feature map is moved from the infrared features $\varphi_i$ to the visible features
$\varphi_v$ in $L=30$ steps and the visible-class probability is recorded on three
MSRS pairs. The probability rises and then saturates near $\varphi_v$, so its
gradient vanishes in the saturated region and the importance map must be accumulated
as an integral gradient rather than a single-point gradient.}\label{fig:csfgrad}

\end{center}

\section{Qualitative comparison across scenes}

\begin{center}

\includegraphics[width=\linewidth]{figures/fig_qual_csf.png}
\captionof{figure}{Qualitative comparison in the style of the classification-saliency
reference: rows are scenes (TNO, RoadScene, MSRS) and columns are methods (VIS, IR,
mean, top-hat, probe, SPGFusion, PIAFusion). Fused panels carry per-panel
CC/SCD/VIF/SSIM values; a red box marks the same target region in every panel, and
the last column shows the zoomed target region of the probe fusion.}\label{fig:qualms}

\end{center}

\section{Component ablation and inference cost}

Table~\ref{tab:ablation} reports the ablation over the individual components of the
rule, and Table~\ref{tab:effic} the inference cost of every method, both measured
under the protocol of the main text.

\begin{center}\small

\captionof{table}{Component ablation. VIF over the three test sets, and downstream
mAP50 and segmentation mIoU. The detection column is the clean-split visible-only
YOLO11n detector of the main text, scored with the same estimator, so the four entries
are the mean, top-hat, probe and ensemble rows of its visible-only block. The
segmentation column is the modality-agnostic segmenter described in the main text. The
control family of Section~3.2 of the main text is evaluated separately in
Table~\ref{tab:segctl}.}\label{tab:ablation}
\begin{tabular*}{\linewidth}{@{\extracolsep{\fill}}l SSS SS@{}}
\toprule
Variant & {VIF-TNO} & {VIF-MSRS} & {VIF-RS} & {mAP50} & {mIoU}\\
\midrule
mean & 0.2684 & 0.2915 & 0.3127 & 0.8273 & 0.1225\\
top-hat prior & 0.3561 & 0.3915 & 0.4241 & \textbf{0.8306} & 0.1270\\
learned probe & 0.3298 & 0.4017 & 0.4227 & 0.8258 & \textbf{0.1270}\\
ensemble & 0.3432 & 0.3956 & 0.4234 & 0.8291 & 0.1269\\
\bottomrule
\end{tabular*}

\end{center}

\begin{table}[!t]\centering\small

\caption{Inference efficiency at $512\times384$, FP32, batch size one, on an
NVIDIA RTX 2080 Ti (lower is better). Latency is averaged over repeated forward passes
after warm-up, and extra peak memory is allocated GPU memory above the pre-call baseline.
SPGFusion parameters count the fusion network only; its memory also includes the DINO
and CLIP feature extractors required at inference. These are the measurements of
the main-text cost table.}
\label{tab:effic}
\begin{tabular*}{\linewidth}{@{\extracolsep{\fill}}l SSS@{}}
\toprule
Method & {Params (M)} & {Time (ms)} & {Extra peak memory (MiB)}\\
\midrule
Ours (mean) & 0.06 & 12.4 & 120.3\\
Ours (top-hat) & 0.06 & 21.2 & 132.1\\
Ours (probe) & 0.07 & 19.1 & 149.3\\
Ours (ensemble) & 0.07 & 41.0 & 188.3\\
IFCNN~\cite{zhang2020} & 0.08 & 7.0 & 242.2\\
DenseFuse~\cite{li2019} & 0.04 & 6.3 & 182.2\\
PIAFusion~\cite{tang2022} & 1.18 & 39.2 & 869.1\\
SPGFusion~\cite{zhang2026} & 75.86 & 538.5 & 1454.8\\
\bottomrule
\end{tabular*}

\end{table}

\section{Qualitative fusion examples}

\begin{center}

\includegraphics[width=0.85\linewidth]{figures/fig_qualitative.png}
\captionof{figure}{Qualitative fusion on a military TNO scene (VIS, IR, Mean, Top-hat,
Probe).}
\label{fig:qual}

\end{center}

\begin{center}

\includegraphics[width=0.85\linewidth]{figures/roadscene.png}
\captionof{figure}{RoadScene qualitative comparison: (a) VIS, (b) IR, (c) top-hat,
(d) probe. The probe trades a little infrared-target fidelity for better
visible-detail retention in this domain-shifted scene.}
\label{fig:road}

\end{center}

\begin{center}

\includegraphics[width=0.92\linewidth]{figures/fig_qual_sota.png}
\captionof{figure}{Qualitative comparison on a TNO scene. SPGFusion and PIAFusion
produce a higher-contrast, higher-fidelity output (higher VIF, Table~III of the main
text), but this pixel-level fidelity does not translate into better downstream
detection under the visible-trained detector reported in the main text.}
\label{fig:qualsota}

\end{center}

\begin{center}

\includegraphics[width=0.85\linewidth]{figures/fig_motivation.png}
\captionof{figure}{Generic saliency (U2-Net) is diffuse and weak, whereas the top-hat
targetness map is structured and target-oriented on a military TNO scene.}
\label{fig:motivation}

\end{center}

\section{Adaptation variants}

Table~\ref{tab:adaptvar} defines the variants of Section~5 of the main text. Every
variant starts from the same reconstruction-trained 24-channel encoder, decoder and
distilled probe, is fine-tuned for $30$ epochs on the 1500 LLVIP training pairs with
the aspect ratio of the native images preserved, and keeps a reconstruction term of
weight $0.1$ in the objective. The pool is the three LLVIP detectors that the main
text names as the tuning pool, namely the visible-trained YOLO11n, the
visible-trained YOLO26n and the modality-agnostic detector. The RT-DETR detector on
LLVIP and all three FLIR detectors are never part of a pool.

\begin{table}[!t]\centering\small

\caption{The adaptation variants. The pool is the LLVIP visible-trained
YOLO11n, the LLVIP visible-trained YOLO26n and the LLVIP modality-agnostic detector.
Round robin back-propagates the loss of one pool member per step and cycles
through the pool, which estimates the gradient of the mean loss; worst case
evaluates the whole pool at each step and back-propagates only the largest loss.
The probe-only variants freeze the encoder and decoder, so the image synthesis is
preserved.}\label{tab:adaptvar}
\resizebox{\linewidth}{!}{%
\begin{tabular}{@{}l l l l@{}}
\toprule
Variant & Tuned & Tuning targets & Objective\\
\midrule
untuned rule & none & none & reconstruction only\\
probe, one & probe & YOLO11n & detector loss, one evaluator\\
probe, pool & probe & full pool & detector loss, round robin over the pool\\
probe, pool, seeds 1 and 2 & probe & full pool & as probe, pool, training seeds 1 and 2\\
round robin & all modules & full pool & detector loss, round robin over the pool\\
worst case & all modules & full pool & detector loss, min-max over the pool\\
detector-tuned & all modules & YOLO11n & detector loss, one evaluator\\
end-to-end & all modules & none & SSIM and MSE agreement with both sources\\
\bottomrule
\end{tabular}}

\end{table}

\section{The fidelity index and its schedule}

The main text defines the reported fidelity index and states the three respects in
which it departs from the reference pixel VIF, namely a Gaussian width of
$2^{s-1}$ at scale $s$ rather than $\text{window}/5$, no decimation between scales,
and a noise variance of $2.0$ held in the units of the input rather than scaled to
them. To measure how much of the reported ordering depends on those choices, we
rescored every fused variant whose outputs we could recover, on the same 30 LLVIP
test images, three ways: with the released implementation, with the reference width
schedule alone, and with the reference construction, which also decimates between
scales and expresses the noise variance in the units of the input.
Table~\ref{tab:vifsched} reports the result. Our own rules and the two retrained
baselines are available at the native detector resolution, while the two published
end-to-end networks survive only at $400\times512$, so every variant in this table is
scored on the sources and the fused image resized to a common $400\times512$ grid.
That is why the values differ from the $512\,$px fidelity benchmark of the main text.

\begin{table}[!t]\centering\footnotesize

\caption{The reported index, each of its three departures from the reference
construction alone, and the reference construction itself, applied to eight fused
variants on the same 30 LLVIP test images, all resized to a common $400\times512$ grid.
\emph{width only} changes the Gaussian width schedule, \emph{decimation only} adds
decimation between scales, and \emph{noise only} expresses the noise variance in the
units of the input, each against the reported configuration. The first, second and last
columns agree with the same three configurations scored by the released
two-dimensional implementation to within $0.001$.
\emph{decimation only} keeps the ranking; \emph{noise only} reverses it at both ends and
collapses the spread.}\label{tab:vifsched}
\begin{tabular*}{\linewidth}{@{\extracolsep{\fill}}l SSSSS@{}} 
\toprule
Method & {reported} & {width only} & {decimation only} & {noise only} & {reference}\\
\midrule
SPGFusion~\cite{zhang2026} & \textbf{0.8043} & \textbf{0.8043} & \textbf{0.8066} & 0.3665 & 0.2410\\
PIAFusion~\cite{tang2022} & 0.6606 & 0.6619 & 0.6555 & \textbf{0.4617} & \textbf{0.3222}\\
IFCNN~\cite{zhang2020} & 0.4850 & 0.4863 & 0.4833 & 0.3941 & 0.3020\\
learned probe & 0.4247 & 0.4103 & 0.4300 & 0.3994 & 0.2969\\
ensemble & 0.4139 & 0.4010 & 0.4180 & 0.4012 & 0.2994\\
top-hat prior & 0.4040 & 0.3924 & 0.4070 & 0.4004 & 0.2992\\
DenseFuse~\cite{li2019} & 0.3700 & 0.3665 & 0.3681 & 0.4026 & 0.3121\\
mean & 0.3410 & 0.3415 & 0.3368 & 0.3998 & 0.3005\\
\bottomrule
\end{tabular*}

\end{table}

Three things follow. The width schedule alone is inconsequential: it leaves the
ordering unchanged and moves the values by at most $0.021$. So does the decimation
between scales, which also leaves the ordering unchanged and moves the values by at
most $0.007$: the reported index and the two constructions that change only its width
schedule or only its decimation all rank SPGFusion first and PIAFusion second. The
third departure is the one that carries the reversal. Changing only the noise
variance, with the widths and the absence of decimation of the reported index, moves
SPGFusion from first to last ($0.8043$ to $0.3665$), moves PIAFusion from second to
first ($0.6606$ to $0.4617$), and collapses the spread across the eight variants from
$0.463$ to $0.095$, which is what the reference construction does as well ($0.081$).
The calibration is therefore responsible for the reversal, and the reference column's
remaining difference from the calibration-only column is a further shift in the
absolute values, not a further change in the ordering. This decomposes the departure
that a three-column rescore leaves open, and it also shows that the reordering is not
a statement about which family of rules wins, because neither construction places our
own rules at the top.

The consequence is stated in the main text. The counterexample holds for the index
the article reports, which is the index under which SPGFusion is the highest-scoring
fused input, and it does not hold for the reference pixel VIF, under which SPGFusion
is the lowest-scoring fused input and therefore agrees with the detection result
rather than contradicting it. The reason is the calibration of the noise variance. The
reported index holds it at the value defined for 8-bit input while the images range
over $[0,1]$, so the noise term is far larger relative to the local variance than in
the reference construction, which leaves the index nearly independent of the assumed
noise level; and the calibration-only column shows that this is the departure that
reverses the ranking, not the width schedule and not the decimation. The absence of
decimation is a construction choice and the calibration is one we did not revisit
before the runs, so the practical reading is narrower than the counterexample alone
suggests: a ranking produced by this index does not transfer to a construction that
calibrates the noise to the input range, and the construction has to be declared with
the numbers. The rescoring script is in the release, under
\texttt{analysis/vif\_reference\_rescore.py}, and the decomposition here is reproduced
by the same script with a single departure changed at a time.

\section{Full reference VIF on every benchmark}

The main-text complete fidelity table now reports both constructions for all fourteen LLVIP
inputs over all 300 test images. Table~\ref{tab:refbench} repeats the comparison for the
eight methods in the three fidelity benchmarks. All three images of each source-family
record are placed on the visible-source grid, capped at 512 px, so that the only
changed variable between columns is the index construction.

The endpoint reversal is identical on all three datasets: SPGFusion is first under the
reported index and last under the reference construction. PIAFusion is the reference
leader on TNO and MSRS; IFCNN leads on RoadScene. This rules out the possibility that
the LLVIP reversal was caused by the particular retained subset or by the
$400\times512$ scoring grid used in the earlier schedule analysis.

\begin{center}\small

\begin{minipage}{\linewidth}
\captionof{table}{Reported and reference VIF on the complete fidelity benchmarks.
The values are means over all pairs; both columns use identical fused outputs and pixels.}
\label{tab:refbench}
\begin{tabular*}{\linewidth}{@{\extracolsep{\fill}}l l SS@{}}
\toprule
Dataset & Method & {reported} & {reference}\\
\midrule
TNO (74) & mean & 0.2683 & 0.2226\\
 & top-hat & 0.3560 & 0.2307\\
 & probe & 0.3297 & 0.2309\\
 & ensemble & 0.3431 & 0.2327\\
 & IFCNN & 0.4440 & 0.2896\\
 & DenseFuse & 0.3697 & 0.2729\\
 & PIAFusion & 0.5485 & \textbf{0.3256}\\
 & SPGFusion & \textbf{0.7855} & 0.1607\\
\addlinespace
MSRS (361) & mean & 0.2914 & 0.2558\\
 & top-hat & 0.3914 & 0.2573\\
 & probe & 0.4015 & 0.2562\\
 & ensemble & 0.3955 & 0.2583\\
 & IFCNN & 0.3757 & 0.2517\\
 & DenseFuse & 0.3004 & 0.2722\\
 & PIAFusion & 0.7507 & \textbf{0.4468}\\
 & SPGFusion & \textbf{0.8314} & 0.2308\\
\addlinespace
RoadScene (221) & mean & 0.3126 & 0.2661\\
 & top-hat & 0.4240 & 0.2617\\
 & probe & 0.4229 & 0.2602\\
 & ensemble & 0.4234 & 0.2637\\
 & IFCNN & 0.4988 & \textbf{0.3082}\\
 & DenseFuse & 0.3828 & 0.3052\\
 & PIAFusion & 0.6315 & 0.2855\\
 & SPGFusion & \textbf{0.8293} & 0.2166\\
\bottomrule
\end{tabular*}

\end{minipage}
\end{center}

\section{Registration sensitivity on LLVIP}

The public LLVIP release contains 3,463 test pairs. We retrieved the source JPEGs for
the same 300 stems used in the detection analysis and estimated a translational
residual with OpenCV ECC on normalized Sobel-edge maps, resized to a 320 px longest
side. The median residual is 2.53 px and the 90th percentile is 4.59 px. We define the
best and worst quartiles by this residual among estimates with ECC at least 0.10; all
300 pairs produced an estimate. The selection never uses labels, fusion scores or
detector results.

Table~\ref{tab:align} gives mAP50 contrasts against the colour-scored visible source.
For the clean-split rows, brackets are observed minima and maxima over four independent
detector runs; the pooled rows use image-level 95 percent paired bootstrap intervals.

\begin{center}\scriptsize

\captionof{table}{Detection sensitivity to source alignment on LLVIP (75 images in each
quartile). Entries are mAP50 differences from the colour-scored visible source.}
\label{tab:align}
\resizebox{\linewidth}{!}{%
\begin{tabular}{@{}l cccc@{}}
\toprule
Evaluator and subset & mean & probe & PIAFusion & SPGFusion\\
\midrule
pooled visible, best & $+.056[-.012,+.111]$ & $+.051[-.018,+.109]$
  & $+.030[-.047,+.098]$ & $-.004[-.086,+.072]$\\
pooled visible, worst & $-.114[-.152,-.075]$ & $-.109[-.143,-.071]$
  & $-.168[-.224,-.109]$ & $-.203[-.263,-.136]$\\
clean visible, best & $+.024[+.011,+.030]$ & $+.032[+.023,+.045]$
  & $+.029[+.002,+.039]$ & $-.021[-.040,+.005]$\\
clean visible, worst & $-.048[-.088,-.027]$ & $-.043[-.052,-.017]$
  & $-.142[-.156,-.108]$ & $-.165[-.190,-.138]$\\
clean agnostic, best & $+.077[+.063,+.102]$ & $+.083[+.066,+.104]$
  & $+.081[+.052,+.117]$ & $+.039[+.007,+.072]$\\
clean agnostic, worst & $+.002[-.044,+.030]$ & $+.013[-.032,+.045]$
  & $+.043[-.012,+.072]$ & $-.004[-.055,+.022]$\\
\bottomrule
\end{tabular}
}

\end{center}

Alignment is therefore a strong conditioning variable rather than a nuisance that can be
assumed away. In the best-aligned clean-split quartile, mean fusion, the top-hat prior
and the probe are above the visible source in every run; in the worst-aligned quartile,
the deficits are largest and resolved. The groups also differ in content: median visible
RMS contrast is 0.085 in the best quartile and 0.110 in the worst, and the corresponding
median annotated box areas are 31,558 and 21,646 pixels. The ECC residual is a
descriptive source-pair estimate, not ground-truth registration metadata, so it cannot
isolate registration from these correlated scene properties. What it does establish is
that the pooled negative result should not be presented as universal across alignment
strata.

\section{Per-seed detection contrasts}

Table~\ref{tab:seeds} prints the individual seed values behind
the main-text seed table. The first four columns are the four independently trained
YOLO11n detectors of the LLVIP benchmark and the next four the four independently
trained YOLO26n detectors, in both cases with seed 0 as the run used for the main
tables and seeds 1 to 3 as the replications.

\begin{table}[!t]\centering\small

\caption{The per-seed contrasts behind the main-text seed table, in mAP50
on the 300 LLVIP test images. Columns 1 to 4 are YOLO11n seeds 0 to 3, columns 5 to
8 are YOLO26n seeds 0 to 3, and the last column is the mean over the four seeds. The
last column is the mean over the four YOLO11n seeds. The ranges quoted in the main text
are the minimum and the maximum of the four per-seed values in that row. The same runs
at the stricter mAP50-95 criterion are summarised in the paragraph below the
table.}\label{tab:seeds}
\resizebox{\linewidth}{!}{%
\begin{tabular}{@{}l cccc cccc S@{}} 
\toprule
Contrast & \multicolumn{4}{c}{YOLO11n} & \multicolumn{4}{c}{YOLO26n} & {mean}\\
\cmidrule(lr){2-5}\cmidrule(lr){6-9}
 & {s0} & {s1} & {s2} & {s3} & {s0} & {s1} & {s2} & {s3} & {}\\
\midrule
SPGFusion $-$ VIS & -0.0908 & -0.0661 & -0.1360 & -0.1520 & -0.0842 & -0.0616 & -0.0353 & -0.0283 & -0.1112\\
mean $-$ VIS & -0.0016 & -0.0142 & -0.0568 & -0.0359 & -0.0245 & -0.0351 & -0.0550 & -0.0381 & -0.0271\\
probe $-$ VIS & -0.0031 & -0.0128 & -0.0431 & -0.0308 & -0.0081 & -0.0209 & -0.0280 & -0.0119 & -0.0224\\
activity-only $-$ VIS & -0.0011 & +0.0004 & -0.0059 & -0.0156 & +0.0100 & +0.0003 & -0.0071 & -0.0292 & -0.0056\\
top-hat $-$ VIS & +0.0017 & -0.0164 & -0.0321 & -0.0370 & -0.0040 & -0.0195 & -0.0214 & -0.0188 & -0.0209\\
probe $-$ mean & -0.0016 & +0.0013 & +0.0137 & +0.0052 & +0.0164 & +0.0142 & +0.0271 & +0.0262 & +0.0047\\
top-hat $-$ mean & +0.0033 & -0.0022 & +0.0247 & -0.0011 & +0.0206 & +0.0156 & +0.0337 & +0.0193 & +0.0062\\
probe $-$ activity-only & -0.0020 & -0.0132 & -0.0372 & -0.0151 & -0.0182 & -0.0213 & -0.0209 & +0.0173 & -0.0169\\
\bottomrule
\end{tabular}}

\end{table}

Four rows of this table carry the reading of the main text. Mean fusion and the
probe are below the unfused visible source at all eight seed and architecture
combinations, while the activity rule exceeds it at one YOLO11n seed and two
YOLO26n seeds and the top-hat prior exceeds it at one YOLO11n seed. Probe minus mean fusion
has a disjoint pair of four-seed ranges across the two architectures, $[-0.0016,
+0.0137]$ against $[+0.0142,+0.0271]$, which is a margin of $0.0005$ before
rounding, and top-hat minus mean fusion overlaps on $[+0.0156,+0.0247]$. The
SPGFusion contrast against the unfused source is negative at every one of the eight
seed and architecture combinations, and the mean fusion and probe contrasts against
the source are negative at every one as well.

The same runs at the stricter mAP50-95 criterion are not as stable, which is why the
main text states that criterion for the reported seed. Over the four YOLO11n seeds the
contrast against the unfused visible source runs from $-0.046$ to $-0.080$ for
SPGFusion, $-0.010$ to $+0.030$ for mean fusion, $-0.003$ to $+0.032$ for the probe,
$+0.008$ to $+0.018$ for the activity control, $-0.002$ to $+0.035$ for the top-hat
prior and $-0.062$ to $+0.007$ for PIAFusion, so the sign of the gain flips at one
seed for mean fusion, the prior and the probe and at two seeds for PIAFusion. Over the
three replicated YOLO26n seeds the activity control, the top-hat prior, the probe and
DenseFuse are ahead of the source at every seed, while the infrared source is behind it
at every seed. No paired interval is computed at this criterion, because the recorded
per-image statistics behind the mAP50 intervals were not retained at both thresholds.

\section{Segmentation results for the control family}

Table~\ref{tab:segctl} reports the control family of Section~3.2 of the main text
under the segmentation evaluator, together with the paired intervals that
Section~4.3 of the main text quotes.

\begin{table}[!t]\centering\small

\caption{Segmentation mIoU on the 361 MSRS-test images for the control
family of Section~3.2 of the main text, with the modality-agnostic segmenter. The
paired intervals resample the 361 images.}\label{tab:segctl}
\begin{tabular*}{\linewidth}{@{\extracolsep{\fill}}l l l@{}} 
\toprule
Variant & {mIoU} & {against mean}\\
\midrule
mean & 0.1225 & \\
activity-only, $S\equiv1$ & 0.1226 & $+0.0002$ $[-0.0353,+0.0358]$\\
top-hat prior & \textbf{0.1270} & $+0.0045$ $[-0.0089,+0.0365]$\\
learned probe & \textbf{0.1270} & $+0.0045$ $[-0.0089,+0.0424]$\\
top-hat, permuted & 0.1228 & \\
probe, permuted & 0.1230 & \\
top-hat, weight-shuffled & 0.1223 & \\
probe, weight-shuffled & 0.1225 & \\
\bottomrule
\end{tabular*}

\end{table}

The two contrasted pairs of the main text are at the top of this table: probe minus
mean is $+0.0045$ with $[-0.0089,+0.0424]$ and top-hat minus activity-only is
$+0.0043$ with $[-0.0182,+0.0480]$, and every interval in the family crosses zero.

\section{Fidelity intervals for the scrambling controls}

Table~\ref{tab:scr} prints the paired intervals behind the statement in
Section~3.2 of the main text that destroying the spatial alignment moves the
reported index back towards mean fusion on every benchmark. Each cell is the mean of
the paired bootstrap distribution of the stated difference and its interval, over
the 74 TNO, 361 MSRS-test and 221 RoadScene images.

\begin{center}\small

\captionof{table}{VIF differences for the scrambling controls, with paired intervals.
Every interval excludes zero, and every value is a difference against the rule
itself, so a positive value means that destroying the alignment costs fidelity.
Benjamini--Hochberg adjusted $q$ values for these comparisons are all below $0.001$
in the recorded analysis.}\label{tab:scr}
\resizebox{\linewidth}{!}{%
\begin{tabular}{@{}l ccc@{}} 
\toprule
Contrast & {TNO} & {MSRS-test} & {RoadScene}\\
\midrule
top-hat $-$ top-hat, permuted & $+0.054$ $[+0.050,+0.058]$ & $+0.098$ $[+0.095,+0.102]$ & $+0.083$ $[+0.080,+0.087]$\\
top-hat $-$ top-hat, weight-shuffled & $+0.083$ $[+0.078,+0.089]$ & $+0.100$ $[+0.097,+0.104]$ & $+0.110$ $[+0.106,+0.114]$\\
probe $-$ probe, permuted & $+0.030$ $[+0.026,+0.034]$ & $+0.107$ $[+0.103,+0.110]$ & $+0.083$ $[+0.080,+0.086]$\\
probe $-$ probe, weight-shuffled & $+0.060$ $[+0.057,+0.064]$ & $+0.108$ $[+0.105,+0.112]$ & $+0.109$ $[+0.105,+0.112]$\\
\bottomrule
\end{tabular}}

\end{center}

\section{Protocol audit of the detection comparison}

Every detection number in the main text is produced by converting each input to a
single channel and stacking it to three, because the fused images are grayscale and
the comparison would otherwise mix conventions. The detectors, however, are trained
on the colour visible JPEGs of the LLVIP subset, so the unfused visible source, which
is the baseline of every comparison in that section, is the one input that could have
been scored in the convention it was trained on. We therefore re-scored the same 300
test images, with the same labels and the same checkpoint, through the library
validator at $640$\,px and a confidence threshold of $0.001$, once in colour and once
in grayscale. Table~\ref{tab:audit} reports the result.

\begin{center}\small

\captionof{table}{The same 300 LLVIP test images, the same checkpoints, and the same
evaluation settings, once with the unfused visible source in the colour convention
the detectors were trained on and once with every input in the grayscale convention
the main tables use. mAP50. Every row is at the native test resolution; the two
published end-to-end networks were re-fused at that resolution rather than upscaled
from $400\times512$.}\label{tab:audit}
\begin{tabular*}{\linewidth}{@{\extracolsep{\fill}}l cc@{}} 
\toprule
Input & {visible-trained} & {modality-agnostic}\\
\midrule
visible source, colour & \textbf{0.8837} & 0.8734\\
visible source, grayscale & 0.8166 & 0.7750\\
infrared source & 0.6968 & \textbf{0.9478}\\
mean fusion & 0.8335 & 0.8879\\
top-hat prior & 0.8332 & 0.8954\\
learned probe & 0.8374 & 0.8950\\
ensemble & 0.8354 & 0.8967\\
IFCNN & 0.7918 & 0.8949\\
DenseFuse & 0.8252 & 0.8823\\
PIAFusion & 0.7627 & 0.9036\\
SPGFusion & 0.7159 & 0.8562\\
\bottomrule
\end{tabular*}

\end{center}

Three consequences follow for the main text. The channel convention alone is worth
$0.067$ mAP50 under the visible-trained detector and $0.098$ under the
modality-agnostic one, and every fused image is grayscale by construction, so the
grayscale visible source is the only row that could have been scored as the detector
was trained. Second, with the colour baseline the statement that no fused input
reaches the unfused visible source becomes stronger: every fused input is below the
colour visible source, PIAFusion and SPGFusion by $0.121$ and $0.168$, and the
largest deficit is now theirs rather than the probe's. Third, the modality reversal
keeps its sign but loses most of its size. Under the modality-agnostic detector every
fused input except SPGFusion is above the colour visible source, by $0.009$ to
$0.030$ rather than by $0.11$ to $0.13$, while SPGFusion stays below it by $0.017$.
The infrared source is still the best input, by $0.044$ over the best fused input,
and its margin over the visible source falls from $0.173$ under the grayscale
convention to $0.074$ under the colour one.

Two further protocol facts belong here. The positive training configuration of the
LLVIP detectors takes its validation split from the
same 300 test images that are scored in the tables, so best-checkpoint selection
optimised for the visible input on exactly those scenes and the absolute mAP50 values
are optimistic. The clean protocol in the final section removes this dependence by
carving 150 validation images out of the 1,500 training pairs; the tables in this
section remain as a pooled sensitivity analysis. The probe input under the
modality-agnostic detector is recorded
as $0.8939$ in one result file and $0.8957$ in another, for the same fused-image
directory and the same checkpoint path. That path no longer exists on the machine
that produced the runs, so we cannot attribute the $0.0018$ difference with certainty
to scoring noise rather than to a rebuilt checkpoint, and we treat contrasts of that
size as unresolved rather than as results.

The table also shows what the choice of metric implementation is worth. It is
produced with the detector library's own validator, and the main text's detection
tables are produced with our own matcher and pooled metric. Our route assigns
predictions to ground truth at each IoU threshold by a maximum-cardinality one-to-one
assignment over the valid pairs (the Hungarian assignment on the complement of the
valid-pair mask), pools the per-image true-positive flags over the 300 images and then
computes the average precision with the library's own function, whereas the validator
matches greedily in confidence order and reads each image through its own loader. Held
to one inference path the two agree to about $0.0001$ on published inputs, so the
differences below belong to the two full pipelines rather than to the matching rule
alone. On the same inputs and
checkpoints the two routes differ by up to $0.015$ mAP50 per row on the sources and
the fused rules of the main tables, and by up to $0.035$ on the two published
end-to-end networks, whose lower detection counts make the two AP rules diverge
further, at $0.0307$ on PIAFusion and $0.0222$ on SPGFusion, which are the rows that
carry the largest deficit the article quotes. On the visible source itself the two
routes differ by $0.0189$ under the modality-agnostic detector and by $0.0123$ under
the visible-trained one. The smallest of these numbers is therefore not a bound on
the others, and we set the descriptive floor of each comparison family from the
disagreement of its own rows: $0.015$ for the fused rules of the main tables where
both routes carry them, $0.019$ for a contrast against the visible source, and
$0.035$ for a contrast involving either published end-to-end network. Every one of
these bounds is larger than the probe-against-mean contrast the main text reports
($-0.0015$ under our route, $+0.0039$ under the library route, a change of sign) and
larger than several other contrasts that the main text treats as resolved. Contrasts
below the bound of their own family are read as unresolved in this article, and the
claims that depend on them are marked as descriptive in the main text.

\paragraph{Initial clean-split validator audit.}
To remove the test-set selection of Table~S10 we retrained both
detectors from the same base weights and the same recipe (YOLO11n, $640$\,px, batch
$16$, seed $0$, $40$ epochs), holding out $150$ of the $1500$ training pairs for
validation and leaving the $300$ test images untouched. The first detector sees the
visible modality only, the second sees both modalities of the same $1350$ training
pairs. Table~\ref{tab:auditsplit} reports this initial eight-input library-validator
audit under the same colour convention. The final supplementary section extends the
independent-test protocol to all three visible-trained architectures and every fused
input under the matched estimator of the main tables.

\begin{center}\small

\captionof{table}{Initial clean-split validator audit, with two detectors retrained on $1350$
training and $150$ validation pairs and the $300$ test images held out, the same
recipe for both, differing only in modality exposure. mAP50, library validator.
Every row is at the native test resolution; the two published end-to-end networks
were re-fused at that resolution rather than upscaled from $400\times512$. The colour
visible source is the best input under the visible-only detector and the infrared
source is the best under the modality-agnostic one; six of the eight fused inputs
fall below the colour visible source under the modality-agnostic detector, and all
eight do under the visible-only one.}\label{tab:auditsplit}
\begin{tabular*}{\linewidth}{@{\extracolsep{\fill}}l cc@{}} 
\toprule
Input & {visible-only, clean split} & {modality-agnostic, clean split}\\
\midrule
visible source, colour & \textbf{0.8620} & 0.8442\\
visible source, grayscale & 0.8294 & 0.7040\\
infrared source & 0.7218 & \textbf{0.9072}\\
mean fusion & 0.8366 & 0.8220\\
top-hat prior & 0.8403 & 0.8352\\
learned probe & 0.8406 & 0.8411\\
ensemble & 0.8418 & 0.8415\\
IFCNN & 0.8312 & 0.8538\\
DenseFuse & 0.8460 & 0.8238\\
PIAFusion & 0.7905 & 0.8734\\
SPGFusion & 0.7333 & 0.8002\\
\bottomrule
\end{tabular*}

\end{center}

Three readings. Under the visible-only detector every fused input is below the
unfused visible source, by $0.016$ to $0.129$, and the published end-to-end networks
are again the weakest. Under the modality-agnostic detector the fused inputs move up
relative to the visible source, but only IFCNN and PIAFusion now exceed it, by
$0.010$ to $0.029$, while mean fusion, the top-hat prior, the probe, the ensemble,
DenseFuse and SPGFusion remain below it by $0.003$ to $0.044$. So the reversal the
main text reports is a comparison between two detector families rather than a
property of the fused inputs: it survives in sign for half of the published methods and
disappears for our own rules, and its size is nearer $0.03$ than the $0.09$ to $0.14$
the grayscale comparison reports. The infrared source remains the best input for the
modality-agnostic detector by $0.034$ over the best fused input.

Two scope statements belong with this table. The corrections of this section are for
the first benchmark only. The second benchmark's visible source is also a grayscale
conversion of a colour input, its detectors were trained with a different budget, and
their weights are no longer available on the machine that produced the runs, so the
two-benchmark phrasing of the first headline is corrected on one benchmark and
reported as measured on the other. And the comparison is pooled over the whole test
set: fusion is proposed for the setting in which the visible modality fails, and this
section does not condition on night, low contrast or occlusion, so the negative is an
average over the benchmark and not a statement about those regimes.

The colour-corrected protocol carries paired intervals for every input of the study, and
Section~13 reports them on the route of the main tables, the matcher and pooled metric
compared above, together with the counts: under the three visible-trained architectures
all fourteen fused inputs are below the colour-scored visible source and forty-one of
the forty-two contrasts carry an interval that excludes zero (all forty-two once
Section~17 repeats the RT-DETR-l family over three runs), and under the
modality-agnostic detector of the main tables ten of the fourteen exceed the source by
$0.023$ to $0.040$. At the reported seed of the clean-split detector none clears its
interval; over four clean-split runs four clear it by $0.022$ to $0.029$. The rows on
which the two routes differ most are the published
end-to-end networks: on the library validator they read $0.7627$ and $0.7159$ against
its own colour visible source of $0.8837$, which is where the $0.121$ and $0.168$ of
Table~\ref{tab:audit} come from, and the $0.035$ between the two routes on those two
rows is the largest in this audit. On the clean split the probe row is the one that
changes sign between the routes, $0.8411$ against a colour source of $0.8442$ on the
library route and a resolved $+0.004$ $[-0.023,+0.029]$ on ours, and we state that
difference rather than resolve it. The interval support the lead claim needs is
therefore reported for the colour-corrected protocol on one route, with the route
difference stated rather than absorbed.

\section{Paired contrasts under every detector of the study}

Section~12 reports the corrected protocol on the estimators it was first rescored with.
This section reports the same contrasts on the estimator of the main tables, for every
input and every detector, so that each table value and the interval beside it come from
one route: per image the predictions are matched to the ground truth by the one-to-one
assignment of \texttt{analysis/j1\_eval2.py}, the true-positive flags are pooled over
the resampled images, and the average precision is the detector library's own function.
Every interval resamples the 300 images 1000 times and applies the same draw to both
arms. Table~\ref{tab:allvis} gives the three visible-trained architectures of the study
and Table~\ref{tab:allagn} the two modality-agnostic ones. A star marks an interval that
excludes zero, and a contrast against the visible source is negative when the fused
input scores below it.

Three readings follow. First, the negative result holds under every visible-trained
architecture. All fourteen fused inputs are below the colour-scored visible source, and
forty-one of the forty-two contrasts carry an interval that excludes zero, the
exception being the activity-only control under RT-DETR-l at $-0.013$
$[-0.030,+0.004]$; Section~17 repeats that architecture over three training runs, where
the control sits below the source in every run and is resolved by their mean. The
resolved deficits run from $0.014$ to $0.152$ mAP50, and the
largest are SPGFusion and IFCNN, which the reported index also ranks first and third of
the fused inputs. The independent-test comparison strengthens this sign: all 42
input-evaluator contrasts are negative. The four-run YOLO11n means span $0.046$ to
$0.136$; the YOLO26n and RT-DETR-l single-run ranges are $0.079$ to $0.223$ and
$0.070$ to $0.246$, with all 14 contrasts surviving their separate BH corrections.
Second, the
modality reversal is resolved under the modality-agnostic detector of the main tables,
where ten of the fourteen fused inputs exceed the colour-scored source by $0.023$ to
$0.040$. The clean split leaves four gains, by $0.022$ to $0.029$ in every run, while
the activity-only control loses in every run. Third, the two published end-to-end networks
are among the weakest fused inputs under both YOLO families and the strongest under the
modality-agnostic detector, so the re-ranking the main text reports is a property of
the evaluator within one estimator as well as across the two estimators compared in
Section~12.

\begin{table}[!t]\centering\small

\caption{Pooled sensitivity audit of every input against the colour-scored visible
source under the three visible-trained architectures, on the estimator of the main
tables. The validation split of these runs overlaps the scored images. Cell: paired
difference at mAP50 and its 95 percent interval over 300 images; a star marks an
interval that excludes zero. YOLO11n and YOLO26n are trained on the 1500 visible
training pairs, RT-DETR-l on the same images with a different budget, and all three
build in the negative result. Every column is one training run; Section~17 repeats the
RT-DETR-l column over two further runs. The independent-test replacement is in the
final section.}\label{tab:allvis}
\resizebox{\linewidth}{!}{%
\begin{tabular}{@{}l ccc@{}}
\toprule
Input & {YOLO11n} & {YOLO26n} & {RT-DETR-l}\\
\midrule
mean fusion              & $-0.054$ $[-0.077,-0.028]$* & $-0.076$ $[-0.104,-0.046]$* & $-0.043$ $[-0.066,-0.018]$*\\
activity-only            & $-0.053$ $[-0.068,-0.037]$* & $-0.041$ $[-0.063,-0.019]$* & $-0.013$ $[-0.030,+0.004]$\\
top-hat prior            & $-0.050$ $[-0.074,-0.026]$* & $-0.056$ $[-0.083,-0.029]$* & $-0.032$ $[-0.057,-0.008]$*\\
learned probe            & $-0.055$ $[-0.079,-0.030]$* & $-0.059$ $[-0.086,-0.031]$* & $-0.038$ $[-0.063,-0.015]$*\\
top-hat, permuted        & $-0.048$ $[-0.071,-0.024]$* & $-0.065$ $[-0.092,-0.036]$* & $-0.033$ $[-0.055,-0.009]$*\\
probe, permuted          & $-0.049$ $[-0.072,-0.024]$* & $-0.070$ $[-0.098,-0.041]$* & $-0.029$ $[-0.051,-0.006]$*\\
top-hat, weight-shuffled & $-0.058$ $[-0.083,-0.032]$* & $-0.071$ $[-0.099,-0.041]$* & $-0.036$ $[-0.059,-0.012]$*\\
probe, weight-shuffled   & $-0.055$ $[-0.079,-0.030]$* & $-0.078$ $[-0.106,-0.048]$* & $-0.039$ $[-0.062,-0.014]$*\\
ensemble                 & $-0.051$ $[-0.075,-0.025]$* & $-0.061$ $[-0.089,-0.033]$* & $-0.031$ $[-0.055,-0.005]$*\\
IFCNN                    & $-0.099$ $[-0.124,-0.075]$* & $-0.089$ $[-0.119,-0.060]$* & $-0.044$ $[-0.068,-0.022]$*\\
DenseFuse                & $-0.059$ $[-0.083,-0.034]$* & $-0.042$ $[-0.072,-0.011]$* & $-0.031$ $[-0.051,-0.009]$*\\
classification saliency  & $-0.057$ $[-0.081,-0.031]$* & $-0.080$ $[-0.108,-0.050]$* & $-0.045$ $[-0.068,-0.020]$*\\
PIAFusion                & $-0.088$ $[-0.119,-0.058]$* & $-0.057$ $[-0.088,-0.028]$* & $-0.054$ $[-0.081,-0.027]$*\\
SPGFusion                & $-0.152$ $[-0.187,-0.117]$* & $-0.143$ $[-0.177,-0.109]$* & $-0.096$ $[-0.124,-0.066]$*\\
visible source, grayscale & $-0.051$ $[-0.065,-0.037]$* & $-0.052$ $[-0.068,-0.036]$* & $-0.036$ $[-0.048,-0.024]$*\\
infrared source          & $-0.174$ $[-0.204,-0.144]$* & $-0.207$ $[-0.242,-0.172]$* & $-0.102$ $[-0.134,-0.073]$*\\
\bottomrule
\end{tabular}}

\end{table}

\begin{table}[!t]\centering\small

\caption{The same contrasts under the two modality-agnostic detectors, the pooled one
of the main tables and the one retrained with the test images held out of validation.
Cell and star as in Table~\ref{tab:allvis}. The reversal the main text reports is
resolved in the first column and not in the second.}\label{tab:allagn}
\resizebox{\linewidth}{!}{%
\begin{tabular}{@{}l cc@{}}
\toprule
Input & {modality-agnostic} & {modality-agnostic, clean split}\\
\midrule
mean fusion              & $+0.036$ $[+0.014,+0.059]$* & $-0.002$ $[-0.029,+0.027]$\\
activity-only            & $+0.001$ $[-0.020,+0.021]$  & $-0.040$ $[-0.064,-0.015]$*\\
top-hat prior            & $+0.037$ $[+0.014,+0.060]$* & $+0.005$ $[-0.021,+0.032]$\\
learned probe            & $+0.040$ $[+0.015,+0.063]$* & $+0.004$ $[-0.023,+0.029]$\\
top-hat, permuted        & $+0.023$ $[-0.001,+0.046]$  & $-0.013$ $[-0.040,+0.016]$\\
probe, permuted          & $+0.023$ $[-0.001,+0.046]$  & $-0.011$ $[-0.039,+0.017]$\\
top-hat, weight-shuffled & $+0.025$ $[+0.002,+0.049]$* & $-0.011$ $[-0.039,+0.015]$\\
probe, weight-shuffled   & $+0.026$ $[+0.003,+0.050]$* & $-0.011$ $[-0.037,+0.016]$\\
ensemble                 & $+0.035$ $[+0.011,+0.060]$* & $+0.006$ $[-0.020,+0.031]$\\
IFCNN                    & $+0.034$ $[+0.009,+0.056]$* & $-0.003$ $[-0.028,+0.021]$\\
DenseFuse                & $+0.030$ $[+0.006,+0.052]$* & $-0.008$ $[-0.037,+0.019]$\\
classification saliency  & $+0.036$ $[+0.014,+0.059]$* & $-0.003$ $[-0.031,+0.025]$\\
PIAFusion                & $+0.030$ $[+0.001,+0.055]$* & $-0.001$ $[-0.026,+0.023]$\\
SPGFusion                & $-0.013$ $[-0.040,+0.011]$  & $-0.045$ $[-0.074,-0.018]$*\\
visible source, grayscale & $-0.100$ $[-0.120,-0.078]$* & $-0.122$ $[-0.146,-0.101]$*\\
infrared source          & $+0.087$ $[+0.060,+0.113]$* & $+0.059$ $[+0.033,+0.084]$*\\
\bottomrule
\end{tabular}}

\end{table}

\section{The pooled negative result conditional on the regime}

Fusion is proposed for the setting in which the visible modality fails, and a pooled
average over a benchmark mixes that setting with the easy one. This section conditions
the pooled sensitivity contrasts on three properties of each scene, each computed from the annotation
and the unfused visible source alone before any contrast is made: the RMS contrast of
that source, the number of annotated pedestrians, and the median annotated box area.
Table~\ref{tab:regime} gives, for every regime and every detector, how many of the
fourteen fused inputs are resolved below the colour-scored visible source, resolved
above it, and unresolved, and Table~\ref{tab:regimelow} gives the per-rule contrasts in
the bottom contrast tercile, the regime the introduction is about. Every interval here
resamples the images of its own regime 1000 times and applies the same draw to both
arms. For each visible-trained architecture, we then apply Benjamini--Hochberg to the
84 two-sided bootstrap $p$ values formed by 14 inputs over six overlapping strata. The
families are separate by architecture, the strata are not treated as independent, and
the multiplicity-adjusted counts are read as an exploratory but more conservative
description than the raw interval counts.

Four readings. First, the pooled average hides a stratification of the benchmark. The
number of deficits that are resolved runs from two of fourteen in the bottom contrast
tercile to fourteen of fourteen in the top one, and the largest deficit in the study,
SPGFusion by $0.291$ mAP50 in the lowest contrast bin, sits at the low end of the
contrast range rather than the high end. Second, that stratification is not a monotone
function of visible contrast, and it is not a regime that the contrast variable
identifies. The fifty images between RMS $0.050$ and $0.067$ have all fourteen deficits
resolved while the bin above them, $0.067$ to $0.083$, has none; and inside the
individual LLVIP sequences the sign of the contrast-deficit association differs from
sequence to sequence, positive in four of the seven sequences with at least thirty test
images and negative in the other three, as Table~\ref{tab:within} shows. Visible
contrast is a property of the scene, so on this benchmark it separates sequences as
much as it separates regimes, and the conditional counts of Table~\ref{tab:regime} are
best read as a stratification of the test set rather than as a dose-response in
illumination. Third, what the split does establish is where the modality reversal is
largest: twelve of the fourteen fused inputs exceed the colour-scored source in the
bottom contrast tercile, by $0.093$ to $0.107$, the lowest two bins carry the largest
reversal of the study, up to $+0.144$, and the top tercile carries none
(Table~\ref{tab:bins}). Fourth, no
stratum contains a fused input that beats the unfused visible source under a
visible-trained detector except the activity-only control under RT-DETR-l at $+0.046$
$[+0.009,+0.092]$, so the pooled conclusion is not an artefact of averaging over
strata: what the average conceals is the size of the effect, not its sign.

\begin{center}\small

\captionof{table}{The fourteen fused inputs against the colour-scored visible source,
by regime, under the estimator of the main tables. Cell: the number resolved below the
source, resolved above it, and unresolved, out of fourteen, after Benjamini--Hochberg
correction at $q<0.05$ over the 84 tests formed by 14 inputs and six overlapping strata
for each architecture. The regimes are defined from the annotation and the unfused
visible source alone, before any contrast is computed, and each interval resamples the
images of that regime 1000 times. The validation split of these runs overlaps the
scored images. Raw interval counts and per-test $q$ values are in
\texttt{regime\_bh\_*.json}; the independent-test counterpart is in the final
section.}\label{tab:regime}
\resizebox{\linewidth}{!}{%
\begin{tabular}{@{}l c ccc@{}}
\toprule
Regime & {n} & {YOLO11n} & {YOLO26n} & {RT-DETR-l}\\
\midrule
low visible contrast, bottom tercile & 100 & 2 / 0 / 12 & 2 / 0 / 12 & 0 / 1 / 13\\
high visible contrast, top tercile   & 100 & 14 / 0 / 0 & 14 / 0 / 0 & 14 / 0 / 0\\
one pedestrian                       & 98  & 0 / 0 / 14 & 0 / 0 / 14 & 1 / 0 / 13\\
three or more pedestrians            & 112 & 14 / 0 / 0 & 14 / 0 / 0 & 5 / 0 / 9\\
small targets                        & 150 & 14 / 0 / 0 & 14 / 0 / 0 & 14 / 0 / 0\\
large targets                        & 150 & 3 / 0 / 11 & 1 / 0 / 13 & 0 / 0 / 14\\
\bottomrule
\end{tabular}}

\end{center}

\begin{table}[!t]\centering\footnotesize

\caption{The same contrasts in the bottom tercile of visible contrast, the
regime the introduction is about, rule by rule. Cell: paired difference at mAP50 and its
95 percent interval over the 100 images of the regime; a star marks $q<0.05$ after the
84-test Benjamini--Hochberg correction for that architecture. Two rules are resolved
below the source under each YOLO family, and the activity-only control is resolved above
it under RT-DETR-l. The two published end-to-end networks stay below the source in all
three point estimates, but their intervals do not all survive correction under
RT-DETR-l.}\label{tab:regimelow}
\resizebox{\linewidth}{!}{%
\begin{tabular}{@{}l ccc@{}}
\toprule
Input & {YOLO11n} & {YOLO26n} & {RT-DETR-l}\\
\midrule
SPGFusion                & $-0.160$ $[-0.258,-0.056]$* & $-0.191$ $[-0.288,-0.081]$* & $-0.077$ $[-0.163,+0.015]$\\
IFCNN                    & $-0.113$ $[-0.184,-0.041]$* & $-0.084$ $[-0.152,-0.007]$* & $-0.008$ $[-0.068,+0.057]$\\
PIAFusion                & $-0.083$ $[-0.173,+0.014]$ & $-0.075$ $[-0.160,+0.029]$ & $-0.038$ $[-0.122,+0.050]$\\
learned probe            & $-0.009$ $[-0.075,+0.066]$ & $-0.058$ $[-0.132,+0.021]$ & $+0.014$ $[-0.056,+0.091]$\\
ensemble                 & $-0.013$ $[-0.077,+0.064]$ & $-0.053$ $[-0.124,+0.028]$ & $+0.019$ $[-0.051,+0.095]$\\
top-hat prior            & $-0.017$ $[-0.081,+0.059]$ & $-0.038$ $[-0.106,+0.039]$ & $+0.011$ $[-0.060,+0.089]$\\
mean fusion              & $-0.023$ $[-0.087,+0.051]$ & $-0.061$ $[-0.132,+0.019]$ & $+0.012$ $[-0.057,+0.088]$\\
classification saliency  & $-0.025$ $[-0.090,+0.048]$ & $-0.061$ $[-0.133,+0.021]$ & $+0.012$ $[-0.058,+0.088]$\\
DenseFuse                & $-0.029$ $[-0.095,+0.047]$ & $-0.013$ $[-0.080,+0.059]$ & $+0.018$ $[-0.046,+0.095]$\\
activity-only            & $-0.025$ $[-0.062,+0.014]$ & $+0.004$ $[-0.045,+0.057]$ & $+0.046$ $[+0.009,+0.092]$*\\
top-hat, permuted        & $-0.024$ $[-0.088,+0.049]$ & $-0.048$ $[-0.118,+0.029]$ & $+0.023$ $[-0.044,+0.097]$\\
probe, permuted          & $-0.023$ $[-0.086,+0.050]$ & $-0.053$ $[-0.124,+0.024]$ & $+0.025$ $[-0.043,+0.100]$\\
top-hat, weight-shuffled & $-0.034$ $[-0.104,+0.043]$ & $-0.052$ $[-0.123,+0.028]$ & $+0.016$ $[-0.052,+0.092]$\\
probe, weight-shuffled   & $-0.028$ $[-0.093,+0.046]$ & $-0.068$ $[-0.140,+0.014]$ & $+0.015$ $[-0.054,+0.091]$\\
\bottomrule
\end{tabular}}

\end{table}

\begin{table}[!t]\centering\small

\caption{The same comparison on a continuous axis: the 300 test images split
into six bins of fifty by the visible source's RMS contrast. Cell: how many of the
fourteen fused inputs are resolved below the colour-scored visible source, resolved
above it, and unresolved, and the paired difference of SPGFusion. The bins are also
sequences, as the last column of Table~\ref{tab:within} shows, so this is a
stratification of the test set rather than a dose-response in illumination.}\label{tab:bins}
\resizebox{\linewidth}{!}{%
\begin{tabular}{@{}c c c c c c@{}}
\toprule
Bin & {RMS contrast} & {YOLO11n below/above/unres.} & {SPGFusion} & {modality-agnostic below/above/unres.} & {SPGFusion}\\
\midrule
1 & 0.050--0.067 & 14/0/0 & $-0.291$ & 0/12/2 & $+0.010$\\
2 & 0.067--0.083 & 0/0/14 & $-0.021$ & 0/13/1 & $+0.059$\\
3 & 0.083--0.118 & 13/0/1 & $-0.128$ & 0/5/9 & $+0.034$\\
4 & 0.124--0.195 & 14/0/0 & $-0.126$ & 4/0/10 & $-0.074$\\
5 & 0.196--0.216 & 12/0/2 & $-0.159$ & 0/0/14 & $-0.036$\\
6 & 0.217--0.299 & 12/0/2 & $-0.156$ & 6/0/8 & $-0.060$\\
\bottomrule
\end{tabular}}

\end{table}

\begin{table}[!t]\centering\small

\caption{The contrast-deficit association inside the seven LLVIP sequences
that contribute at least thirty test images. Each sequence is split at its own median
visible contrast, and the two halves are scored separately against the colour-scored
visible source. Columns 4 and 5 give the paired difference of SPGFusion in the
lower-contrast and higher-contrast halves; the last column is the difference between
the two for SPGFusion under the modality-agnostic detector. The sign of the association
is not the same in every sequence, which is why Section~14 reads the conditioning as a
stratification rather than a regime effect.}\label{tab:within}
\resizebox{\linewidth}{!}{%
\begin{tabular}{@{}c c c c c c c@{}}
\toprule
Sequence & {n} & {contrast median} & {SPG, low half} & {SPG, high half} & {mean, low half} & {SPG low $-$ high, agnostic}\\
\midrule
190 & 78 & 0.206 & $-0.089$ & $-0.214$ & $-0.019$ & $-0.012$\\
200 & 31 & 0.083 & $+0.268$ & $-0.069$ & $+0.281$ & $+0.200$\\
210 & 32 & 0.092 & $-0.224$ & $-0.135$ & $-0.015$ & $-0.061$\\
220 & 38 & 0.175 & $-0.144$ & $-0.241$ & $-0.170$ & $+0.076$\\
230 & 46 & 0.066 & $-0.196$ & $-0.040$ & $-0.085$ & $-0.083$\\
240 & 31 & 0.071 & $-0.350$ & $-0.371$ & $-0.217$ & $-0.118$\\
260 & 44 & 0.203 & $-0.317$ & $+0.003$ & $-0.074$ & $+0.082$\\
\bottomrule
\end{tabular}}

\end{table}

\section{The same stratification in a second task}

The detection comparison is one task with one evaluator family. This section repeats
the contrast split on semantic segmentation, the second task of the article, to see
whether the stratification is a property of the detector or of the comparison. The
361 MSRS-test images are split into terciles of the unfused visible source's RMS
contrast, and every input is scored with the segmenter of Section~4.3 of the main text;
the paired difference against the unfused visible source is computed within each
tercile with image-resampled intervals. Two conventions matter here and are stated in
both places. The mIoU used in this section averages over the classes present in each
label, the usual convention; the main text's segmentation table was produced with the
stricter convention of its own script, which counts a class absent from both the
prediction and the label as an IoU of zero, so its absolute values are lower by
construction. That difference is not cosmetic: under the all-classes convention of the
main table the infrared source is the best input, and under the present-classes
convention used here the unfused visible source is the best input everywhere, which is
one more instance of the article's own point that a ranking is a property of the metric
as much as of the input.

Table~\ref{tab:segstrat} gives the result. The stratification of the detection deficit
reproduces in sign and in order: the fused inputs are indistinguishable from the
unfused visible source in the bottom contrast tercile, where the largest deficit is
$0.006$ mIoU and no interval excludes zero, they lose $0.013$ to $0.026$ in the middle
tercile and $0.032$ to $0.054$ in the top tercile, where every interval excludes zero.
Under this task and this metric the targetness rules are the least penalised of the
fused inputs in every stratum, as they are under detection, and the infrared source is
never the best input once the convention is stated. The cross-task reading is
therefore that the size of the fusion penalty tracks how much the visible source has to
offer, in a second task and a second benchmark, while the sign of the source ranking
itself remains a property of the metric.

\begin{table}[!t]\centering\small

\caption{Segmentation mIoU on the 361 MSRS-test images, split into terciles of
the visible source's RMS contrast (thresholds $0.135$ and $0.201$). Cell: mean mIoU and
the paired difference against the unfused visible source with its 95 percent interval;
a star marks an interval that excludes zero. mIoU is averaged over the classes present
in each label, which is the convention of this section and not the stricter convention
of the main text's segmentation table.}\label{tab:segstrat}
\resizebox{\linewidth}{!}{%
\begin{tabular}{@{}l ccc@{}}
\toprule
Input & {low contrast, bottom tercile} & {middle tercile} & {high contrast, top tercile}\\
\midrule
visible source & \textbf{0.2944} & \textbf{0.3147} & \textbf{0.3460}\\
infrared source & $0.2695$ ($-0.025$ $[-0.039,-0.011]$*) & $0.2731$ ($-0.042$ $[-0.064,-0.018]$*) & $0.2582$ ($-0.088$ $[-0.114,-0.062]$*)\\
mean fusion & $0.2900$ ($-0.004$ $[-0.018,+0.009]$) & $0.2898$ ($-0.025$ $[-0.047,-0.004]$*) & $0.2924$ ($-0.054$ $[-0.074,-0.034]$*)\\
activity-only & $0.2924$ ($-0.002$ $[-0.015,+0.010]$) & $0.2930$ ($-0.022$ $[-0.042,-0.001]$*) & $0.2952$ ($-0.051$ $[-0.067,-0.035]$*)\\
top-hat prior & $0.2888$ ($-0.006$ $[-0.018,+0.007]$) & $0.2975$ ($-0.017$ $[-0.037,+0.003]$) & $0.3124$ ($-0.033$ $[-0.050,-0.017]$*)\\
learned probe & $0.2890$ ($-0.006$ $[-0.017,+0.007]$) & $0.3016$ ($-0.013$ $[-0.032,+0.006]$) & $0.3137$ ($-0.032$ $[-0.048,-0.017]$*)\\
top-hat, permuted & $0.2902$ ($-0.004$ $[-0.018,+0.009]$) & $0.2955$ ($-0.019$ $[-0.041,+0.002]$) & $0.2978$ ($-0.048$ $[-0.068,-0.030]$*)\\
probe, permuted & $0.2891$ ($-0.005$ $[-0.019,+0.008]$) & $0.2958$ ($-0.019$ $[-0.040,+0.002]$) & $0.3001$ ($-0.046$ $[-0.065,-0.028]$*)\\
top-hat, weight-shuffled & $0.2899$ ($-0.004$ $[-0.018,+0.009]$) & $0.2889$ ($-0.026$ $[-0.048,-0.004]$*) & $0.2927$ ($-0.053$ $[-0.074,-0.034]$*)\\
probe, weight-shuffled & $0.2892$ ($-0.005$ $[-0.018,+0.008]$) & $0.2909$ ($-0.024$ $[-0.045,-0.002]$*) & $0.2952$ ($-0.051$ $[-0.071,-0.032]$*)\\
\bottomrule
\end{tabular}}

\end{table}

\section{The clean-split modality reversal across seeds}

Section~13 reports the modality-agnostic detector retrained on a clean split at its
reported seed, where no fused input clears its interval. Three further seeds of the same
detector were trained with the same recipe on the same 1350 training pairs and the same
150 held-out validation pairs, and Table~\ref{tab:seedsplit} gives each seed's paired
contrast against the colour-scored visible source together with a two-level bootstrap of
the mean over the four seeds, in which seeds are drawn with replacement and the images
of each drawn seed are resampled 1000 times.

With the seed dimension included, four of the fourteen fused inputs exceed the source by
an interval that excludes zero: the top-hat prior by $0.025$, the ensemble by $0.024$,
the probe by $0.022$ and PIAFusion by $0.029$. The activity-only control is resolved
below the source by $0.022$, and the remaining nine inputs are indistinguishable from
it. The corrected modality reversal therefore survives the clean split for four rules,
at a size of $0.022$ to $0.029$ mAP50, which is the resolution bound the main tables
report, and the reported seed alone understates it: that seed is the only one of the
four in which no fused input clears its interval, and the three added seeds are the
reason the reversal is resolved at all. The second finding of the article therefore rests
on the pooled modality-agnostic detector and on the seed-averaged clean split, not on a
single retraining run, and its size after both corrections is about $0.025$ mAP50
rather than the $0.09$ to $0.14$ the grayscale comparison reports.

\begin{table}[!t]\centering\small

\caption{The clean-split modality-agnostic detector over four seeds. Columns 1
to 4 are the paired difference against the colour-scored visible source for seed 0, the
seed reported in Section~13, and the three seeds added here, in mAP50. The last column
is the mean over the four seeds with its 95 percent interval from a two-level bootstrap
that resamples seeds and then images within each drawn seed; a star marks an interval
that excludes zero.}\label{tab:seedsplit}
\resizebox{\linewidth}{!}{%
\begin{tabular}{@{}l SSSS l@{}} 
\toprule
Input & {seed 0} & {seed 1} & {seed 2} & {seed 3} & {mean over seeds}\\
\midrule
mean fusion              & +0.000 & +0.023 & +0.031 & +0.037 & $+0.023$ $[-0.001,+0.042]$\\
activity-only            & -0.038 & -0.013 & -0.020 & -0.014 & $-0.022$ $[-0.042,-0.002]$*\\
top-hat prior            & +0.006 & +0.030 & +0.032 & +0.036 & $+0.025$ $[+0.004,+0.044]$*\\
learned probe            & +0.005 & +0.027 & +0.025 & +0.033 & $+0.022$ $[+0.003,+0.041]$*\\
top-hat, permuted        & -0.011 & +0.016 & +0.025 & +0.023 & $+0.013$ $[-0.010,+0.033]$\\
probe, permuted          & -0.009 & +0.018 & +0.026 & +0.026 & $+0.015$ $[-0.008,+0.035]$\\
top-hat, weight-shuffled & -0.009 & +0.021 & +0.024 & +0.024 & $+0.014$ $[-0.008,+0.034]$\\
probe, weight-shuffled   & -0.008 & +0.016 & +0.025 & +0.026 & $+0.014$ $[-0.009,+0.034]$\\
ensemble                 & +0.007 & +0.026 & +0.030 & +0.035 & $+0.024$ $[+0.004,+0.042]$*\\
IFCNN                    & -0.002 & +0.017 & +0.022 & +0.024 & $+0.015$ $[-0.004,+0.034]$\\
DenseFuse                & -0.007 & +0.015 & +0.023 & +0.021 & $+0.013$ $[-0.009,+0.032]$\\
classification saliency  & -0.001 & +0.024 & +0.030 & +0.037 & $+0.022$ $[-0.002,+0.042]$\\
PIAFusion                & +0.000 & +0.028 & +0.044 & +0.042 & $+0.029$ $[+0.003,+0.051]$*\\
SPGFusion                & -0.045 & -0.004 & +0.003 & +0.008 & $-0.010$ $[-0.041,+0.015]$\\
\bottomrule
\end{tabular}}

\end{table}

\section{Repetition of the pooled RT-DETR-l family over three runs}

Section~13 reports one training run of each visible-trained architecture, and under
RT-DETR-l one of the forty-two contrasts, the activity-only control, keeps an interval
that includes zero. That architecture has the smallest deficits of the three, so the
reading of the negative result under it turns on whether those deficits are larger than
its own run-to-run spread. Two further runs of the same recipe were therefore trained on
the same 1500 visible training pairs and scored the same way, per image and against that
run's own colour-scored visible source. Table~\ref{tab:rtdetrseeds} gives the three runs
and the two-level bootstrap of their mean, in which runs are drawn with replacement and
the images of each drawn run are resampled 1000 times.

Two readings. First, the deficits of this architecture are stable across runs and stay
the mildest of the three: SPGFusion loses $0.096$, $0.074$ and $0.104$ mAP50 in the three
runs, against $0.152$ under YOLO11n on the same images, and its mean over runs is
$-0.091$ $[-0.115,-0.067]$. No run puts a fused input above its own visible source, and
the two published end-to-end networks are the two weakest fused inputs in every run as
well. The architecture therefore joins the other two in sign while remaining the
smallest effect, which is what Section~13 reads it as. Second, the control row that
carries the exception is below the source in all three runs, by $0.013$, $0.018$ and
$0.011$ mAP50, and the mean over runs excludes zero at $-0.014$ $[-0.027,-0.001]$. Each
run on its own stays unresolved, so what the three runs change is the resolution of that
one contrast rather than its direction: with the run dimension included, all forty-two
contrasts of Section~13 are resolved below the colour-scored visible source. The
qualifier that belongs with this row is that the colour-scored visible source itself
moves by $0.020$ mAP50 between the three runs ($0.907$, $0.915$ and $0.895$), which is
larger than the control's deficit in each run, so the paired design is what makes the
comparison, and a resolved interval on a deficit of $0.014$ is a statement at the edge of
the resolution this benchmark offers.

\begin{center}\small

\captionof{table}{The pooled RT-DETR-l family over three training runs, on the same 300
LLVIP test images and against each run's own colour-scored visible source. Validation
overlaps the scored images; the independent-test result is in the final section.
Columns 1 to 3 are
the paired difference at mAP50 for the reported run and the two runs added here, and the
last column is their mean with its 95 percent interval from a two-level bootstrap that
resamples runs and then images within each drawn run; a star marks an interval that
excludes zero. The reported run's own intervals, which agree with
Table~\ref{tab:allvis}, are not repeated here.}\label{tab:rtdetrseeds}
\resizebox{\linewidth}{!}{%
\begin{tabular}{@{}l SSS l@{}}
\toprule
Input & {run 0} & {run 1} & {run 2} & {mean over runs}\\
\midrule
mean fusion              & -0.044 & -0.028 & -0.040 & $-0.037$ $[-0.056,-0.019]$*\\
activity-only            & -0.013 & -0.018 & -0.011 & $-0.014$ $[-0.027,-0.001]$*\\
top-hat prior            & -0.032 & -0.022 & -0.041 & $-0.032$ $[-0.050,-0.013]$*\\
learned probe            & -0.038 & -0.020 & -0.035 & $-0.031$ $[-0.050,-0.012]$*\\
top-hat, permuted        & -0.033 & -0.032 & -0.027 & $-0.031$ $[-0.047,-0.014]$*\\
probe, permuted          & -0.029 & -0.028 & -0.027 & $-0.028$ $[-0.044,-0.011]$*\\
top-hat, weight-shuffled & -0.037 & -0.032 & -0.026 & $-0.032$ $[-0.049,-0.014]$*\\
probe, weight-shuffled   & -0.040 & -0.031 & -0.026 & $-0.032$ $[-0.051,-0.014]$*\\
ensemble                 & -0.031 & -0.015 & -0.038 & $-0.028$ $[-0.048,-0.008]$*\\
IFCNN                    & -0.044 & -0.043 & -0.047 & $-0.045$ $[-0.061,-0.029]$*\\
DenseFuse                & -0.031 & -0.013 & -0.020 & $-0.022$ $[-0.039,-0.005]$*\\
classification saliency  & -0.045 & -0.026 & -0.042 & $-0.038$ $[-0.058,-0.018]$*\\
PIAFusion                & -0.054 & -0.062 & -0.068 & $-0.061$ $[-0.082,-0.044]$*\\
SPGFusion                & -0.096 & -0.074 & -0.104 & $-0.091$ $[-0.115,-0.067]$*\\
\bottomrule
\end{tabular}}

\end{center}

\section{Method details and algorithm listing}
\label{sec:lmethod}

The compact main text keeps only what is needed to read the results, and this section
carries the configuration behind Section III of the main text. The encoder is a stride-1
convolution stack with reflection padding, Leaky-ReLU (slope $0.2$) and $5\times5$
kernels; its channel widths are $1,16,24,24,24$ and its output is the 24-channel
feature volume used by the fusion rule. The decoder mirrors it with widths
$24,16,16,16,1$ and outputs a single channel in $[-1,1]$ through a final $\tanh$. Both
are trained once on $128\times128$ patches of the MSRS training set with
$\mathcal{L}_{\mathrm{rec}}=\lambda\,\mathrm{MSE}+(1-\mathrm{SSIM})$, $\lambda=25$, and
are frozen afterwards. The learned probe is two $3\times3$ convolutions
($24\rightarrow16$ followed by Leaky-ReLU, then $16\rightarrow1$ without activation)
trained by distillation against the top-hat prior with the encoder frozen. The
morphological prior uses a $9\times9$ rectangular structuring element and a Gaussian of
$\sigma=3$, with the per-image min--max normalisation applied last. All fusion results
use $\tau=0.3$. Algorithm~1 is the complete per-pair procedure.

\begin{center}
\fbox{\parbox{0.94\linewidth}{\small
\textbf{Algorithm 1} \ \textit{Target-aware semantic saliency fusion (per pair)}\\[1pt]
\textbf{Input:} visible $I_v$, infrared $I_r$; encoder $E$; decoder $D$; targetness
(morphological top-hat prior or learned probe $\Phi$); $\tau=0.3$.\\
\textbf{Output:} fused image $I_f$.\\
1: $F_v\gets E(I_v)$,\; $F_r\gets E(I_r)$.\\
2: for $s\in\{v,r\}$: $S_s\gets$ top-hat prior or $\Phi(F_s)$;\; $A_s\gets\sum_c|F_s|$.\\
3: $d_s\gets S_s\odot A_s$;\; $m\gets\max(\max_x d_v,\max_x d_r)$;\;
$w_s\gets \exp\big(d_s/(m\tau)\big)/\sum_{s'}\exp\big(d_{s'}/(m\tau)\big)$.\\
4: $\hat{F}\gets\sum_s w_s\odot F_s$;\; $I_f\gets D(\hat{F})$.
}}
\end{center}

\section{Detection under the reported protocol at native resolution}
\label{sec:lnative}

Table~\ref{tab:lnative} reports the reported protocol, that is, with every input scored
in the single-channel convention that treats all rows alike, under the primary
visible-trained detector and under the modality-agnostic detector of the main tables.
Under the visible-trained detector the method the reported index ranks first is the
worst input: SPGFusion is $0.101$ below the unfused visible source with a 95 percent
interval of $[-0.135,-0.066]$ over the 300 images, IFCNN is $0.048$
$[-0.074,-0.019]$ below it and PIAFusion $0.037$ $[-0.068,-0.009]$ below it, while
every other fused input, the eight control variants included, lies within $0.008$
mAP50 of the source and no interval separates it. Under the modality-agnostic detector
the same table reverses: the infrared source becomes the best input of all at $0.9415$
against $0.7561$ for the visible source, and ten of the fourteen fused inputs exceed
the visible source. The colour-corrected protocol described in the main text is the one
uses for its negative result, because it removes the penalty that the single-channel
convention imposes on the visible source.

\begin{table}[!t]\centering\small

\caption{Detection on the 300-image LLVIP test set at native resolution,
reported protocol. The left block scores every input with the visible-trained YOLO11n
and the right block with the modality-agnostic detector; the two further
visible-trained architectures are in Table~\ref{tab:allvis} of Sec.~13. All fused
images come from one checkpoint set and all values from one metric implementation.
mAP50 and mAP50-95.}\label{tab:lnative}
\setlength{\tabcolsep}{4pt}
\resizebox{\linewidth}{!}{%
\begin{tabular}{@{}l SSSS@{}}
\toprule
& \multicolumn{2}{c}{visible-trained} & \multicolumn{2}{c}{modality-agnostic}\\
\cmidrule(lr){2-3}\cmidrule(lr){4-5}
Input & {mAP50} & {mAP50-95} & {mAP50} & {mAP50-95}\\
\midrule
VIS source & 0.8289 & 0.4500 & 0.7561 & 0.4014\\
IR source & 0.7080 & 0.3911 & \textbf{0.9415} & \textbf{0.6168}\\
mean fused & 0.8273 & 0.4799 & 0.8912 & 0.5365\\
activity-only fused & 0.8277 & 0.4676 & 0.8555 & 0.5064\\
top-hat fused & 0.8306 & 0.4844 & 0.8912 & 0.5398\\
probe fused & 0.8258 & 0.4820 & 0.8939 & 0.5481\\
ensemble fused & 0.8291 & 0.4820 & 0.8953 & 0.5461\\
top-hat, permuted & 0.8331 & 0.4758 & 0.8776 & 0.5306\\
probe, permuted & 0.8320 & 0.4772 & 0.8775 & 0.5355\\
top-hat, weight-shuffled & 0.8237 & 0.4755 & 0.8795 & 0.5304\\
probe, weight-shuffled & 0.8259 & 0.4787 & 0.8798 & 0.5307\\
IFCNN fused & 0.7816 & 0.4469 & 0.8880 & 0.5470\\
DenseFuse fused & 0.8283 & 0.4829 & 0.8868 & 0.5329\\
CSF fused & 0.8244 & 0.4795 & 0.8909 & 0.5331\\
PIAFusion fused & 0.7929 & 0.4561 & 0.8845 & 0.5563\\
SPGFusion fused & 0.7296 & 0.4006 & 0.8412 & 0.5207\\
\bottomrule
\end{tabular}}

\end{table}

\section{The second detection benchmark}
\label{sec:lflir}

The second benchmark is FLIR-ADAS aligned, in the registered release with COCO
annotations, from which we select the 1500 training and 300 test pairs that contain at
least one person box, at native $640\times512$, so that the benchmark has the same
shape as the first: 1500 training images, 300 test images, one class. The first
benchmark is drawn the same way from the LLVIP test split, so both are subsets of the
published splits selected to contain a target, and their absolute mAP50 values are not
comparable with numbers reported on the full splits. Two detectors are trained on the
second benchmark with the same recipe, one on the visible modality only and one on both,
and a third visible-trained detector differs from the first only in architecture. Ten
fused inputs are scored by all three, and the two published networks are re-fused on
that benchmark at native resolution.

\begin{table}[!t]\centering\small

\caption{Replication on FLIR-ADAS aligned, 1500 training and 300 test pairs,
single class person, native $640\times512$. Three detectors are trained on this
benchmark with the same recipe, two of them on the visible modality alone and differing
in architecture. mAP50; the paired intervals are in the text of the main text and in
Table~\ref{tab:lcross}. Best in bold.}\label{tab:lflir}
\resizebox{\linewidth}{!}{%
\begin{tabular}{@{}l ccc@{}}
\toprule
Input & {YOLO11n} & {YOLO26n} & {agnostic}\\
\midrule
VIS source & \textbf{0.4686} & \textbf{0.3769} & 0.3950\\
IR source & 0.3382 & 0.2955 & \textbf{0.7630}\\
mean fused & 0.3834 & 0.3424 & 0.6511\\
activity-only fused & 0.4018 & 0.3541 & 0.6162\\
top-hat fused & 0.4156 & 0.3571 & 0.6305\\
probe fused & 0.4188 & 0.3686 & 0.6255\\
ensemble fused & 0.4156 & 0.3666 & 0.6296\\
IFCNN fused & 0.4062 & 0.3449 & 0.6217\\
DenseFuse fused & 0.4131 & 0.3636 & 0.6486\\
CSF fused & 0.3813 & 0.3215 & 0.6463\\
PIAFusion fused & 0.3773 & 0.3544 & 0.5528\\
SPGFusion fused & 0.3758 & 0.3459 & 0.5359\\
\bottomrule
\end{tabular}}

\end{table}

Two readings. Under both visible-trained architectures the unfused visible source is
the best of the twelve inputs and every fused input lies below it. Under YOLO11n the
tabulated mean differences (fusion minus VIS) run from $-0.0498$ for the probe to
$-0.0928$ for SPGFusion; paired bootstrap estimates for these two endpoints are
$-0.0515$ $[-0.0840,-0.0161]$ and $-0.0928$ $[-0.1251,-0.0627]$, so both exclude zero.
Under YOLO26n the point-estimate differences run from $-0.0084$ for the probe to
$-0.0554$ for the classification-saliency rule; intervals were not computed for this
evaluator. The two visible-trained architectures do not agree on the internal
ordering: PIAFusion is tenth of twelve under YOLO11n and sixth under YOLO26n, and the
spread between the best and the worst fused input is $0.043$ against $0.047$. With the
modality-agnostic detector the ranking moves again, the infrared source becomes the
best input by a wide margin, every fused input gains on the visible source by $+0.14$
to $+0.26$, and the targetness rules that beat mean fusion under the visible-trained
detector become worse than it.

\begin{table}[!t]\centering\small

\caption{The same five contrasts under all four combinations of detection
benchmark and evaluator. A star marks a paired interval that excludes zero; each cell
is the mean of the bootstrap distribution of that contrast on the same estimator as its
interval, which can differ from the difference of the two pooled table values by up to
$0.003$. Since the two metric routes disagree by up to $0.015$ on the fused rules of
the main tables, the starred cells whose size is near that bound are read as descriptive
rather than as resolved. Every contrast changes sign at least once.}\label{tab:lcross}
\begin{tabular*}{\linewidth}{@{\extracolsep{\fill}}l cccc@{}}
\toprule
Contrast & {LLVIP, vis} & {LLVIP, agn} & {FLIR, vis} & {FLIR, agn}\\
\midrule
SPGFusion $-$ VIS & $-0.0927$* & $+0.0924$* & $-0.0928$* & $+0.1404$*\\
mean $-$ VIS & $-0.0023$ & $+0.1330$* & $-0.0862$* & $+0.2556$*\\
probe $-$ mean & $-0.0017$ & $+0.0031$ & $+0.0347$* & $-0.0265$*\\
probe $-$ activity-only & $-0.0027$ & $+0.0388$* & $+0.0172$* & $+0.0089$\\
top-hat $-$ mean & $+0.0033$ & $+0.0003$ & $+0.0316$* & $-0.0211$*\\
\bottomrule
\end{tabular*}

\end{table}

\section{Adaptation: per-evaluator results and image drift}
\label{sec:ladapt}

Table~\ref{tab:ladapt} gives mAP50 under every evaluator for the untuned rule and the
four adapted variants, and Table~\ref{tab:ldrift} gives how far the fused image moves.
The untuned row is scored inside the adaptation harness, which is why its LLVIP
modality-agnostic entry reads $0.8957$ here where the main tables read $0.8939$ from
the main pass; every contrast of this section is computed within the harness pass.

\begin{table}[!t]\centering\small

\caption{mAP50 under every evaluator for the untuned rule and four adapted
variants. The last three columns are a different benchmark, and the RT-DETR column is a
held-out LLVIP evaluator. The variants are tuned against a pool drawn from the LLVIP
YOLO11n, YOLO26n and modality-agnostic detectors, so the RT-DETR column and the three
FLIR columns are held out and marked with a dagger.}\label{tab:ladapt}
\resizebox{\linewidth}{!}{%
\begin{tabular}{@{}l cccc ccc@{}}
\toprule
& \multicolumn{4}{c}{LLVIP} & \multicolumn{3}{c}{FLIR}\\
\cmidrule(lr){2-5}\cmidrule(lr){6-8}
Variant & {11n} & {26n} & {RT-DETR$^\dagger$} & {agn} & {11n$^\dagger$} & {26n$^\dagger$} & {agn$^\dagger$}\\
\midrule
untuned rule & 0.8258 & 0.8128 & 0.8651 & 0.8957 & 0.4189 & 0.3686 & 0.6255\\
probe, one & 0.8103 & 0.8535 & 0.8641 & 0.8423 & 0.4184 & 0.3597 & 0.6049\\
probe, pool & 0.8348 & 0.8528 & 0.8775 & 0.8840 & 0.4164 & 0.3459 & 0.6065\\
round robin & 0.8261 & 0.8118 & 0.8831 & 0.8891 & 0.1530 & 0.2153 & 0.3122\\
worst case & 0.8184 & 0.7901 & 0.8553 & 0.8283 & 0.1754 & 0.2086 & 0.2085\\
\bottomrule
\end{tabular}}

\end{table}

\begin{sloppypar}
Adapting only the probe moves the rule between evaluators without raising the average:
the pooled variant gains $+0.0389$ on the visible-trained YOLO26n with a paired interval
of $[+0.0208,+0.0598]$ and loses $0.0114$ $[-0.0243,-0.0007]$ on the LLVIP
modality-agnostic detector, $0.0224$ $[-0.0352,-0.0093]$ on the held-out FLIR YOLO26n
detector and $0.0188$ $[-0.0347,-0.0029]$ on the held-out FLIR modality-agnostic
detector, leaving the remaining three evaluators unchanged within their intervals. The
same sign pattern holds for two further seeds, $+0.0363$ and $+0.0427$ on YOLO26n and
$-0.0161$ and $-0.0236$ on the modality-agnostic detector. Adapting the whole network is
a different regime: round robin and worst case leave the LLVIP evaluators roughly
intact but collapse on the second benchmark, from $0.4189$ to $0.1530$ on the
visible-trained detector and from $0.6255$ to $0.3122$ on the modality-agnostic one.
\end{sloppypar}

\begin{table}[!t]\centering\small

\caption{How far the fused image moves when the rule adapts. Drift is the
SSIM against the untuned output; the last column is the standard deviation of the output
on the second benchmark, next to $0.1568$ for the untuned rule.}\label{tab:ldrift}
\resizebox{\linewidth}{!}{%
\begin{tabular}{@{}l cccc c@{}}
\toprule
& \multicolumn{2}{c}{LLVIP} & \multicolumn{3}{c}{FLIR}\\
\cmidrule(lr){2-3}\cmidrule(lr){4-6}
Variant & {drift SSIM} & {abs. diff.} & {drift SSIM} & {abs. diff.} & {output s.d.}\\
\midrule
probe, pool & 0.9402 & 0.0317 & 0.9238 & 0.0331 & 0.1716\\
probe, pool, seed 1 & 0.9279 & 0.0369 & 0.9319 & 0.0408 & 0.1857\\
probe, pool, seed 2 & 0.9410 & 0.0326 & 0.9397 & 0.0371 & 0.1833\\
round robin & 0.8169 & 0.0631 & 0.7660 & 0.1535 & 0.1172\\
worst case & 0.7884 & 0.0795 & 0.6376 & 0.2601 & 0.0967\\
\bottomrule
\end{tabular}}

\end{table}

\section{Independent-test replication across major LLVIP evaluators}
\label{sec:cleanr2}

The initial clean-split validator audit used the two YOLO11n evaluators and
eight fused outputs. To make the independent test set the primary basis of every major
LLVIP evaluator comparison, we fixed one split before retraining the missing
evaluators. The split contains $1{,}350$ training, $150$ validation and $300$ test
images with $4{,}479$, $419$ and $744$ pedestrian boxes. The stem sets are pairwise
disjoint, and \texttt{clean\_split\_manifest.json} records all 1,800 selected stems.
The 300 test images are exactly those used by the fixed fused-image archive.

All three visible-only evaluators start from their official pretrained checkpoints.
Their runs use 40 epochs at $640$\,px and seed 0. YOLO11n and YOLO26n use batch size 16;
RT-DETR-l uses batch size 4 to fit the 11 GB GPU. YOLO11n is repeated for seeds 1 to 3.
The modality-agnostic control is a four-seed YOLO11n trained on both registered
modalities of the same 1,350 pairs. The library selects each \texttt{best.pt} using
validation fitness on the disjoint 150 images. No test image enters optimisation,
checkpoint selection, threshold selection or calibration. All predictions use
Ultralytics 8.4.189 at confidence $0.001$ and input size $640$.

Every one of the 14 fused inputs and both source images is scored under every
evaluator. The primary contrast is paired over the same 300 test images against each
checkpoint's colour-scored visible source. For each single-run clean evaluator we apply
Benjamini--Hochberg separately to its 14 fused-input tests. The frozen per-image
statistics, intervals, $q$ values, training manifests and checkpoint hashes are in
\texttt{results/clean\_r2/}.

\begin{table}[!t]\centering\small

\caption{Independent-test contrasts for all 14 fused inputs. The visible-only and
modality-agnostic YOLO11n columns give the mean over four runs and the full observed
run range. The YOLO26n and RT-DETR-l columns give the observed paired difference and
the image-bootstrap 95 percent interval; a star marks $q<0.05$ in a separate 14-test
Benjamini--Hochberg family for each one-run evaluator. The header gives the visible
source mAP50.}\label{tab:cleanr2}
\resizebox{\linewidth}{!}{%
\begin{tabular}{@{}l cccc@{}}
\toprule
Input & {\shortstack[c]{YOLO11n VIS\\4 runs (0.8657)}} &
{\shortstack[c]{YOLO26n VIS\\1 run (0.8122)}} &
{\shortstack[c]{RT-DETR-l VIS\\1 run (0.8095)}} &
{\shortstack[c]{YOLO11n agn.\\4 runs (0.8660)}}\\
\midrule
mean fusion & $-0.060$ $[-0.086,-0.029]$ & \shortstack[l]{$-0.147$ $[-0.188,-0.107]$\\$q<0.002$*} & \shortstack[l]{$-0.129$ $[-0.166,-0.090]$\\$q<0.002$*} & $+0.023$ $[-0.000,+0.038]$\\
activity-only & $-0.046$ $[-0.062,-0.026]$ & \shortstack[l]{$-0.087$ $[-0.119,-0.055]$\\$q<0.002$*} & \shortstack[l]{$-0.070$ $[-0.100,-0.044]$\\$q<0.002$*} & $-0.021$ $[-0.038,-0.013]$\\
top-hat prior & $-0.048$ $[-0.074,-0.024]$ & \shortstack[l]{$-0.148$ $[-0.190,-0.104]$\\$q<0.002$*} & \shortstack[l]{$-0.123$ $[-0.162,-0.086]$\\$q<0.002$*} & $+0.026$ $[+0.006,+0.036]$\\
learned probe & $-0.054$ $[-0.082,-0.028]$ & \shortstack[l]{$-0.142$ $[-0.185,-0.101]$\\$q<0.002$*} & \shortstack[l]{$-0.119$ $[-0.159,-0.083]$\\$q<0.002$*} & $+0.023$ $[+0.005,+0.034]$\\
top-hat, permuted & $-0.058$ $[-0.088,-0.039]$ & \shortstack[l]{$-0.130$ $[-0.170,-0.092]$\\$q<0.002$*} & \shortstack[l]{$-0.110$ $[-0.147,-0.073]$\\$q<0.002$*} & $+0.014$ $[-0.011,+0.026]$\\
probe, permuted & $-0.060$ $[-0.095,-0.034]$ & \shortstack[l]{$-0.126$ $[-0.163,-0.090]$\\$q<0.002$*} & \shortstack[l]{$-0.108$ $[-0.145,-0.072]$\\$q<0.002$*} & $+0.015$ $[-0.010,+0.026]$\\
top-hat, w-shuffled & $-0.062$ $[-0.088,-0.037]$ & \shortstack[l]{$-0.143$ $[-0.183,-0.103]$\\$q<0.002$*} & \shortstack[l]{$-0.124$ $[-0.163,-0.087]$\\$q<0.002$*} & $+0.015$ $[-0.010,+0.025]$\\
probe, w-shuffled & $-0.065$ $[-0.096,-0.038]$ & \shortstack[l]{$-0.144$ $[-0.186,-0.104]$\\$q<0.002$*} & \shortstack[l]{$-0.121$ $[-0.159,-0.084]$\\$q<0.002$*} & $+0.015$ $[-0.008,+0.026]$\\
ensemble & $-0.050$ $[-0.080,-0.023]$ & \shortstack[l]{$-0.145$ $[-0.189,-0.102]$\\$q<0.002$*} & \shortstack[l]{$-0.110$ $[-0.149,-0.074]$\\$q<0.002$*} & $+0.025$ $[+0.008,+0.035]$\\
IFCNN & $-0.076$ $[-0.091,-0.050]$ & \shortstack[l]{$-0.092$ $[-0.132,-0.052]$\\$q<0.002$*} & \shortstack[l]{$-0.095$ $[-0.131,-0.059]$\\$q<0.002$*} & $+0.016$ $[-0.001,+0.024]$\\
DenseFuse & $-0.052$ $[-0.078,-0.029]$ & \shortstack[l]{$-0.079$ $[-0.118,-0.041]$\\$q<0.002$*} & \shortstack[l]{$-0.088$ $[-0.126,-0.053]$\\$q<0.002$*} & $+0.013$ $[-0.007,+0.023]$\\
classif. saliency & $-0.061$ $[-0.089,-0.030]$ & \shortstack[l]{$-0.149$ $[-0.189,-0.109]$\\$q<0.002$*} & \shortstack[l]{$-0.132$ $[-0.170,-0.094]$\\$q<0.002$*} & $+0.023$ $[-0.002,+0.037]$\\
PIAFusion & $-0.104$ $[-0.123,-0.075]$ & \shortstack[l]{$-0.198$ $[-0.244,-0.151]$\\$q<0.002$*} & \shortstack[l]{$-0.207$ $[-0.252,-0.163]$\\$q<0.002$*} & $+0.029$ $[+0.001,+0.045]$\\
SPGFusion & $-0.136$ $[-0.146,-0.126]$ & \shortstack[l]{$-0.223$ $[-0.269,-0.177]$\\$q<0.002$*} & \shortstack[l]{$-0.246$ $[-0.290,-0.204]$\\$q<0.002$*} & $-0.009$ $[-0.044,+0.009]$\\
\bottomrule
\end{tabular}}

\end{table}

\begin{sloppypar}
Table~\ref{tab:cleanr2} completes the evaluator coverage requested for the first
benchmark. Under YOLO26n and RT-DETR-l every one of the 14 deficits is significant in
its own BH family, from $0.079$ to $0.246$ mAP50 across the two evaluators. The
YOLO11n four-run means span $0.046$ to $0.136$ and remain negative in every observed
run. The ranking is also partly architecture-dependent: the Spearman association of
the 14 contrasts is $0.45$ between YOLO11n and YOLO26n, $0.60$ between YOLO11n and
RT-DETR-l, and $0.94$ between YOLO26n and RT-DETR-l. These are 14-rule descriptive
associations, not a population estimate. Under the modality-agnostic evaluator, two
inputs are below the visible source and twelve have positive point estimates, of which
four are positive in every run and all four are resolved by the two-level interval
(Table~\ref{tab:seedsplit}).
\end{sloppypar}

\begin{center}\small

\captionof{table}{The same six overlapping scene regimes under the independent-test
evaluators. Each cell gives the number resolved below by BH at $q<0.05$, resolved above,
and unresolved, out of fourteen. The family is 14 inputs by six strata for each
evaluator. The YOLO11n column uses its reported clean seed; YOLO26n and RT-DETR-l each
have one clean run.}\label{tab:cleanregime}
\resizebox{\linewidth}{!}{%
\begin{tabular}{@{}l cccccc@{}}
\toprule
Evaluator & {\shortstack[c]{low contrast\\$n=100$}} &
{\shortstack[c]{high contrast\\$n=100$}} &
{\shortstack[c]{one person\\$n=98$}} &
{\shortstack[c]{three or more\\$n=112$}} &
{\shortstack[c]{small targets\\$n=150$}} &
{\shortstack[c]{large targets\\$n=150$}}\\
\midrule
YOLO11n & 1 / 0 / 13 & 1 / 0 / 13 & 0 / 0 / 14 & 3 / 0 / 11 & 2 / 0 / 12 & 0 / 0 / 14\\
YOLO26n & 13 / 0 / 1 & 14 / 0 / 0 & 11 / 0 / 3 & 14 / 0 / 0 & 14 / 0 / 0 & 14 / 0 / 0\\
RT-DETR-l & 2 / 0 / 12 & 14 / 0 / 0 & 2 / 0 / 12 & 14 / 0 / 0 & 14 / 0 / 0 & 5 / 0 / 9\\
\bottomrule
\end{tabular}}

\end{center}

\begin{sloppypar}
The clean stratification is not a single architecture-invariant partition. YOLO26n
resolves most strata, whereas the clean YOLO11n seed resolves only one to three inputs
in the contrast and size strata, and RT-DETR-l sits between them. The common pattern is
narrower: the negative aggregate result holds under every visible-trained clean
evaluator, but which subgroup supplies statistical resolution changes with the
detector. This is why the main text retains the regime counts as exploratory evidence
rather than as a universal mechanism.
\end{sloppypar}

\begin{thebibliography}{10}
\expandafter\ifx\csname url\endcsname\relax
  \def\url#1{\texttt{#1}}\fi
\expandafter\ifx\csname urlprefix\endcsname\relax\def\urlprefix{URL }\fi
\expandafter\ifx\csname href\endcsname\relax
  \def\href#1#2{#2} \def\path#1{#1}\fi

\bibitem{ma2019}
J.~Ma, Y.~Ma, C.~Li, {Infrared and visible image fusion methods and
  applications: A survey}, Information Fusion 45 (2019) 153--178.
\newblock \href {https://doi.org/10.1016/j.inffus.2018.02.004}
  {\path{doi:10.1016/j.inffus.2018.02.004}}.

\bibitem{gajjar2026visual}
T.~Gajjar and M.~Goyani, {Infrared--visible image fusion for robust visual
  perception: Methods, benchmarks, evaluation and open challenges}, The Visual
  Computer 42(11) (2026).
\newblock \href {https://doi.org/10.1007/s00371-026-04695-5}
  {\path{doi:10.1007/s00371-026-04695-5}}.

\bibitem{zhu2017}
P.~Zhu, X.~Ma, Z.~Huang, {Fusion of infrared-visible images using improved
  multi-scale top-hat transform and suitable fusion rules}, Infrared Physics \&
  Technology 81 (2017) 282--295.
\newblock \href {https://doi.org/10.1016/j.infrared.2017.01.013}
  {\path{doi:10.1016/j.infrared.2017.01.013}}.

\bibitem{zhang2021}
S.~Zhang, X.~Li, X.~Zhang, S.~Zhang, {Infrared and visible image fusion based
  on saliency detection and two-scale transform decomposition}, Infrared
  Physics \& Technology 114 (2021) 103626.
\newblock \href {https://doi.org/10.1016/j.infrared.2020.103626}
  {\path{doi:10.1016/j.infrared.2020.103626}}.

\bibitem{yin2013}
H.~Yin, S.~Li, L.~Fang, {Simultaneous image fusion and super-resolution using
  sparse representation}, Information Fusion 14~(3) (2013) 229--240.
\newblock \href {https://doi.org/10.1016/j.inffus.2012.01.008}
  {\path{doi:10.1016/j.inffus.2012.01.008}}.

\bibitem{nejati2015}
M.~Nejati, S.~Samavi, S.~Shirani, {Multi-focus image fusion using
  dictionary-based sparse representation}, Information Fusion 25 (2015) 72--84.
\newblock \href {https://doi.org/10.1016/j.inffus.2014.10.004}
  {\path{doi:10.1016/j.inffus.2014.10.004}}.

\bibitem{yin2015}
H.~Yin, {Sparse representation with learned multiscale dictionary for image
  fusion}, Neurocomputing 148 (2015) 600--610.
\newblock \href {https://doi.org/10.1016/j.neucom.2014.07.003}
  {\path{doi:10.1016/j.neucom.2014.07.003}}.

\bibitem{wang2014}
L.~Wang, B.~Li, L.-f. Tian, {Multi-modal medical image fusion using the
  inter-scale and intra-scale dependencies between image shift-invariant
  shearlet coefficients}, Information Fusion 19 (2014) 20--28.
\newblock \href {https://doi.org/10.1016/j.inffus.2012.03.002}
  {\path{doi:10.1016/j.inffus.2012.03.002}}.

\bibitem{zhang2016}
Q.~Zhang, X.~Maldague, {An adaptive fusion approach for infrared and visible
  images based on NSCT and compressed sensing}, Infrared Physics \& Technology
  74 (2016) 11--20.
\newblock \href {https://doi.org/10.1016/j.infrared.2015.11.003}
  {\path{doi:10.1016/j.infrared.2015.11.003}}.

\bibitem{wei2017}
C.~Wei, B.~Zhou, W.~Guo, {A three scale image transformation for infrared and
  visible image fusion}, 2017 20th International Conference on Information
  Fusion (Fusion) (2017) 1--6\newblock \href {https://doi.org/10.23919/icif.2017.8009748}
  {\path{doi:10.23919/icif.2017.8009748}}.

\bibitem{fu2009}
M.-Y. FU, C.~ZHAO, {FUSION OF INFRARED AND VISIBLE IMAGES BASED ON THE SECOND
  GENERATION CURVELET TRANSFORM}, JOURNAL OF INFRARED AND MILLIMETER WAVES
  28~(4) (2009) 254--258.
\newblock \href {https://doi.org/10.3724/sp.j.1010.2009.00254}
  {\path{doi:10.3724/sp.j.1010.2009.00254}}.

\bibitem{ma2021}
X.~Ma, Z.~Wang, S.~Hu, {Multi-focus image fusion based on multi-scale sparse
  representation}, Journal of Visual Communication and Image Representation 81
  (2021) 103328.
\newblock \href {https://doi.org/10.1016/j.jvcir.2021.103328}
  {\path{doi:10.1016/j.jvcir.2021.103328}}.

\bibitem{chen2007}
H.~Chen, P.~K. Varshney, {A human perception inspired quality metric for image
  fusion based on regional information}, Information Fusion 8~(2) (2007)
  193--207.
\newblock \href {https://doi.org/10.1016/j.inffus.2005.10.001}
  {\path{doi:10.1016/j.inffus.2005.10.001}}.

\bibitem{piella}
G.~Piella, H.~Heijmans, {A new quality metric for image fusion}, in: Proceedings
  of the 2003 International Conference on Image Processing (ICIP), vol.~3, 2003,
  pp.~173--176.
\newblock \href {https://doi.org/10.1109/ICIP.2003.1247209}
  {\path{doi:10.1109/ICIP.2003.1247209}}.

\bibitem{ma20191}
J.~Ma, W.~Yu, P.~Liang, C.~Li, J.~Jiang, {FusionGAN: A generative adversarial
  network for infrared and visible image fusion}, Information Fusion 48 (2019)
  11--26.
\newblock \href {https://doi.org/10.1016/j.inffus.2018.09.004}
  {\path{doi:10.1016/j.inffus.2018.09.004}}.

\bibitem{hou2021}
J.~Hou, D.~Zhang, W.~Wu, J.~Ma, H.~Zhou, {A Generative Adversarial Network for
  Infrared and Visible Image Fusion Based on Semantic Segmentation}, Entropy
  23~(3) (2021) 376.
\newblock \href {https://doi.org/10.3390/e23030376}
  {\path{doi:10.3390/e23030376}}.

\bibitem{yi2021}
S.~Yi, J.~Li, X.~Yuan, {DFPGAN: Dual fusion path generative adversarial network
  for infrared and visible image fusion}, Infrared Physics \& Technology 119
  (2021) 103947.
\newblock \href {https://doi.org/10.1016/j.infrared.2021.103947}
  {\path{doi:10.1016/j.infrared.2021.103947}}.

\bibitem{song2022}
A.~Song, H.~Duan, H.~Pei, L.~Ding, {Triple-discriminator generative adversarial
  network for infrared and visible image fusion}, Neurocomputing 483 (2022)
  183--194.
\newblock \href {https://doi.org/10.1016/j.neucom.2022.02.025}
  {\path{doi:10.1016/j.neucom.2022.02.025}}.

\bibitem{xu2022}
H.~Xu, J.~Ma, J.~Jiang, X.~Guo, H.~Ling, {U2Fusion: A Unified Unsupervised
  Image Fusion Network}, IEEE Transactions on Pattern Analysis and Machine
  Intelligence 44~(1) (2022) 502--518.
\newblock \href {https://doi.org/10.1109/tpami.2020.3012548}
  {\path{doi:10.1109/tpami.2020.3012548}}.

\bibitem{li2021}
H.~Li, X.-J. Wu, J.~Kittler, {RFN-Nest: An end-to-end residual fusion network
  for infrared and visible images}, Information Fusion 73 (2021) 72--86.
\newblock \href {https://doi.org/10.1016/j.inffus.2021.02.023}
  {\path{doi:10.1016/j.inffus.2021.02.023}}.

\bibitem{li2019}
H.~Li, X.-J. Wu, {DenseFuse: A Fusion Approach to Infrared and Visible Images},
  IEEE Transactions on Image Processing 28~(5) (2019) 2614--2623.
\newblock \href {https://doi.org/10.1109/tip.2018.2887342}
  {\path{doi:10.1109/tip.2018.2887342}}.

\bibitem{zhang2020}
Y.~Zhang, Y.~Liu, P.~Sun, H.~Yan, X.~Zhao, L.~Zhang, {IFCNN: A general image
  fusion framework based on convolutional neural network}, Information Fusion
  54 (2020) 99--118.
\newblock \href {https://doi.org/10.1016/j.inffus.2019.07.011}
  {\path{doi:10.1016/j.inffus.2019.07.011}}.

\bibitem{wang2024}
X.~Wang, L.~Fang, J.~Zhao, Z.~Pan, H.~Li, Y.~Li, {UUD-Fusion: An unsupervised
  universal image fusion approach via generative diffusion model}, Computer
  Vision and Image Understanding 249 (2024) 104218.
\newblock \href {https://doi.org/10.1016/j.cviu.2024.104218}
  {\path{doi:10.1016/j.cviu.2024.104218}}.

\bibitem{liu2023}
R.~Liu, H.~Wang, S.~Du, Y.~Liu, {A Visible and Infrared Image Fusion Framework
  Based on Dual-Path Encoder-Decoder and Multi-Scale Discrete Wavelet
  Transform}, 2023 IEEE International Conference on Image Processing (ICIP)
  (2023) 1995--1999\newblock \href {https://doi.org/10.1109/icip49359.2023.10222863}
  {\path{doi:10.1109/icip49359.2023.10222863}}.

\bibitem{hao2022}
S.~Hao, T.~He, B.~An, X.~Ma, H.~Wen, F.~Wang, {VDFEFuse: A novel fusion
  approach to infrared and visible images}, Infrared Physics \& Technology 121
  (2022) 104048.
\newblock \href {https://doi.org/10.1016/j.infrared.2022.104048}
  {\path{doi:10.1016/j.infrared.2022.104048}}.

\bibitem{yang2023}
X.~Yang, H.~Huo, R.~Wang, C.~Li, X.~Liu, J.~Li, {DGLT-Fusion: A decoupled
  global–local infrared and visible image fusion transformer}, Infrared
  Physics \& Technology 128 (2023) 104522.
\newblock \href {https://doi.org/10.1016/j.infrared.2022.104522}
  {\path{doi:10.1016/j.infrared.2022.104522}}.

\bibitem{an2024}
R.~An, G.~Liu, Y.~Qian, M.~Xing, H.~Tang, {MRASFusion: A multi-scale residual
  attention infrared and visible image fusion network based on semantic
  segmentation guidance}, Infrared Physics \& Technology 139 (2024) 105343.
\newblock \href {https://doi.org/10.1016/j.infrared.2024.105343}
  {\path{doi:10.1016/j.infrared.2024.105343}}.

\bibitem{yi2024}
X.~Yi, L.~Tang, H.~Zhang, H.~Xu, J.~Ma, {Diff-IF: Multi-modality image fusion
  via diffusion model with fusion knowledge prior}, Information Fusion 110
  (2024) 102450.
\newblock \href {https://doi.org/10.1016/j.inffus.2024.102450}
  {\path{doi:10.1016/j.inffus.2024.102450}}.

\bibitem{yang2025}
B.~Yang, Z.~Jiang, D.~Pan, H.~Yu, G.~Gui, W.~Gui, {LFDT-Fusion: A latent
  feature-guided diffusion Transformer model for general image fusion},
  Information Fusion 113 (2025) 102639.
\newblock \href {https://doi.org/10.1016/j.inffus.2024.102639}
  {\path{doi:10.1016/j.inffus.2024.102639}}.

\bibitem{tang2022}
L.~Tang, J.~Yuan, H.~Zhang, X.~Jiang, J.~Ma, {PIAFusion: A progressive infrared
  and visible image fusion network based on illumination aware}, Information
  Fusion 83-84 (2022) 79--92.
\newblock \href {https://doi.org/10.1016/j.inffus.2022.03.007}
  {\path{doi:10.1016/j.inffus.2022.03.007}}.

\bibitem{zhang2026}
H.~Zhang, S.~Yao, J.~Ma, J.~Jiang, Y.~Zhang, and H.~Zhou, {SPGFusion: Semantic
  prior guided infrared and visible image fusion via pretrained vision models},
  Information Fusion 125 (2026) 103433.
\newblock \href {https://doi.org/10.1016/j.inffus.2025.103433}
  {\path{doi:10.1016/j.inffus.2025.103433}}.

\bibitem{zhou2026}
Y.~Zhou, Q.~Yang, K.~Ren, Q.~Chen, {MTFusion: A semantic segmentation-guided
  infrared and visible image fusion network based on an attention-enhanced
  hybrid Mamba-Transformer architecture}, Infrared Physics \& Technology 153
  (2026) 106325.
\newblock \href {https://doi.org/10.1016/j.infrared.2025.106325}
  {\path{doi:10.1016/j.infrared.2025.106325}}.

\bibitem{bianco2020}
S.~Bianco, M.~Buzzelli, G.~Ciocca, R.~Schettini, {Neural architecture search
  for image saliency fusion}, Information Fusion 57 (2020) 89--101.
\newblock \href {https://doi.org/10.1016/j.inffus.2019.12.007}
  {\path{doi:10.1016/j.inffus.2019.12.007}}.

\bibitem{khan2025}
H.~Khan, M.~T. Usman, J.~Koo, {Bilateral Feature Fusion with hexagonal
  attention for robust saliency detection under uncertain environments},
  Information Fusion 121 (2025) 103165.
\newblock \href {https://doi.org/10.1016/j.inffus.2025.103165}
  {\path{doi:10.1016/j.inffus.2025.103165}}.

\bibitem{wang2025}
Q.~Wang, Z.~Li, S.~Zhang, Y.~Luo, W.~Chen, T.~Wang, N.~Chi, Q.~Dai,
  {SMAE-Fusion: Integrating saliency-aware masked autoencoder with hybrid
  attention transformer for infrared–visible image fusion}, Information
  Fusion 117 (2025) 102841.
\newblock \href {https://doi.org/10.1016/j.inffus.2024.102841}
  {\path{doi:10.1016/j.inffus.2024.102841}}.

\bibitem{xu2021}
H.~Xu, H.~Zhang, J.~Ma, {Classification Saliency-Based Rule for Visible and
  Infrared Image Fusion}, IEEE Transactions on Computational Imaging 7 (2021)
  824--836.
\newblock \href {https://doi.org/10.1109/tci.2021.3100986}
  {\path{doi:10.1109/tci.2021.3100986}}.

\bibitem{sheikh2006}
H.~R. Sheikh, A.~C. Bovik, {Image information and visual quality}, IEEE
  Transactions on Image Processing 15~(2) (2006) 430--444.
\newblock \href {https://doi.org/10.1109/TIP.2005.859378}
  {\path{doi:10.1109/TIP.2005.859378}}.

\bibitem{qin2020}
X.~Qin, Z.~Zhang, C.~Huang, M.~Dehghan, O.~R. Zaiane, M.~Jagersand, {U2-Net:
  Going deeper with nested U-structure for salient object detection}, Pattern
  Recognition 106 (2020) 107404.
\newblock \href {https://doi.org/10.1016/j.patcog.2020.107404}
  {\path{doi:10.1016/j.patcog.2020.107404}}.

\bibitem{zhu2025}
J.~Zhu, X.~Qin, A.~Elsaddik, {DC-Net: Divide-and-conquer for salient object
  detection}, Pattern Recognition 157 (2025) 110903.
\newblock \href {https://doi.org/10.1016/j.patcog.2024.110903}
  {\path{doi:10.1016/j.patcog.2024.110903}}.

\bibitem{kirillov2023sam}
A.~Kirillov, E.~Mintun, N.~Ravi, J.~Shlens, J.~Johnson, {Segment Anything}, in:
  Proceedings of the IEEE/CVF International Conference on Computer Vision
  (ICCV), 2023, pp. 4015--4026.
\newblock \href {https://doi.org/10.1109/ICCV51070.2023.00371}
  {\path{doi:10.1109/ICCV51070.2023.00371}}.

\bibitem{guo2024}
J.~Guo, Z.~Huang, Y.~Tao, {Multispectral Pedestrian Detection Based on
  Prior-Saliency Attention and Image Fusion}, Electronics 13~(9) (2024) 1770.
\newblock \href {https://doi.org/10.3390/electronics13091770}
  {\path{doi:10.3390/electronics13091770}}.

\bibitem{petilli2022}
E.~Petilli, {Monocular thermal imaging for pedestrian detection and ranging},
  Optical Engineering 62~(03) (2022) 031212.
\newblock \href {https://doi.org/10.1117/1.OE.62.3.031212}
  {\path{doi:10.1117/1.OE.62.3.031212}}.

\bibitem{li2010}
J.~Li, W.~Gong, W.~Li, X.~Liu, {Robust pedestrian detection in thermal infrared
  imagery using the wavelet transform}, Infrared Physics \& Technology 53~(4)
  (2010) 267--273.
\newblock \href {https://doi.org/10.1016/j.infrared.2010.03.005}
  {\path{doi:10.1016/j.infrared.2010.03.005}}.

\bibitem{redmon2016yolo}
J.~Redmon, S.~Divvala, R.~Girshick, A.~Farhadi, {You Only Look Once: Unified,
  Real-Time Object Detection}, in: Proceedings of the IEEE Conference on
  Computer Vision and Pattern Recognition (CVPR), 2016, pp. 779--788.
\newblock \href {https://doi.org/10.1109/CVPR.2016.91}
  {\path{doi:10.1109/CVPR.2016.91}}.

\bibitem{jocher2024yolo}
G.~Jocher, J.~Qiu, {Ultralytics YOLO}, software release (2024).
\newblock \href {https://github.com/ultralytics/ultralytics}
  {\path{https://github.com/ultralytics/ultralytics}}.

\bibitem{lv2024rtdetr}
W.~Lv, S.~Xu, Y.~Zhao, G.~Wang, J.~Wei, C.~Cui, Y.~Du, Q.~Dang, Y.~Liu,
  {DETRs Beat YOLOs on Real-time Object Detection}, in: Proceedings of the
  IEEE/CVF Conference on Computer Vision and Pattern Recognition (CVPR), 2024,
  pp. 16965--16974.
\newblock \href {https://doi.org/10.1109/CVPR52733.2024.01605}
  {\path{doi:10.1109/CVPR52733.2024.01605}}.

\bibitem{ronneberger2015}
O.~Ronneberger, P.~Fischer, T.~Brox, {U-Net: Convolutional Networks for
  Biomedical Image Segmentation}, Lecture Notes in Computer Science 9351 (2015)
  234--241\newblock \href {https://doi.org/10.1007/978-3-319-24574-4_28}
  {\path{doi:10.1007/978-3-319-24574-4_28}}.

\bibitem{flir_adas}
Teledyne FLIR, {FLIR Thermal Dataset for Algorithm Training}, dataset (2018).
\newblock \href {https://www.flir.com/oem/adas/adas-dataset-form/}
  {\path{https://www.flir.com/oem/adas/adas-dataset-form/}}.

\end{thebibliography}


\end{document}
